\chapter{MMU}
\label{chap:mmu}

\begin{introduction}[本章内容]
  \item 39 位虚拟地址与三级页表
  \item 硬件如何遍历页表、块映射
  \item 描述符格式与内存属性（MAIR）
  \item 清 BSS、启动页表（恒等映射加高地址映射）、TCR/MAIR/SCTLR、跳到高地址
  \item 多核启动：两道闸门
  \item 两阶段的物理内存初始化：\code{struct page}、\code{mem\_map} 与伙伴系统
  \item 最终内核页表与用户页表
  \item 第一个进程、\code{fork} 与地址空间切换；TLB；缺页与按需分配
  \item \code{mmap()}：VMA、按需调页、共享页与设备映射
  \item MMIO、总线地址与 DMA
  \item 三个缓存与 TLB 故障
  \item 实验分支：引用计数与写时复制
\end{introduction}

第~\ref{chap:boot} 章结束时，CPU 已经处在 EL1，但 MMU 还关着。本章先讲 AArch64 怎样翻译地址，再沿着启动顺序看页表怎样一步步建起来：\file{entry.S} 的启动页表、副核唤醒、物理页分配器、\code{main()} 里的最终内核页表，然后是每个进程的用户页表和 \code{mmap()}。最后几节转向页表的另一面——它不只决定地址，还决定 CPU 怎样访问一块内存：设备寄存器、DMA 缓冲区，以及三个只在真机上出现的缓存与 TLB 故障。

\section{AArch64 的地址翻译}

\subsection{地址空间划分}

本内核使用 4\,KiB 页、39 位虚拟地址。\reg{TCR\_EL1} 中 \code{T0SZ}$=$\code{T1SZ}$=25$，即每个地址空间有 $2^{64-25}=2^{39}$ 字节（512\,GiB）。AArch64 用虚拟地址的高位选择页表基址寄存器：

\begin{itemize}
  \item 高位全 0：\code{0x0000000000000000}～\code{0x0000007fffffffff}，由 \reg{TTBR0\_EL1} 翻译，用于\textbf{用户空间}，每个进程一张；
  \item 高位全 1：\code{0xffffff8000000000}～\code{0xffffffffffffffff}，由 \reg{TTBR1\_EL1} 翻译，用于\textbf{内核空间}，所有进程共享一张；
  \item 中间的地址不合法，访问即产生翻译错误。
\end{itemize}

两个基址寄存器分开带来一个很大的好处：进程切换时只需改写 \reg{TTBR0\_EL1}，内核映射保持不变，也不用像 xv6-riscv 那样在每个用户页表里复制一份内核映射和 trampoline 页。

\subsection{三级页表}

39 位地址用 4\,KiB 粒度时，从第 1 级开始遍历，共三级（没有第 0 级）。每级页表占一页，含 512 个 8 字节表项，用虚拟地址中的 9 位索引（图~\ref{fig:va}）：

\begin{figure}[htbp]
  \centering
  \begin{tikzpicture}[x=3.3mm,y=8mm]
    \foreach \x/\w/\t/\c in {0/6/{高位: 全0或全1}/gray!15, 6/9/{L1 索引 VA[38:30]}/structurecolor!15,
        15/9/{L2 索引 VA[29:21]}/structurecolor!25, 24/9/{L3 索引 VA[20:12]}/structurecolor!35,
        33/12/{页内偏移 VA[11:0]}/gray!10}{
      \draw[fill=\c] (\x,0) rectangle ++(\w,1);
      \node[font=\scriptsize, text width=\w*3.1mm, align=center] at (\x+\w/2,0.5) {\t};
    }
    \node[lab, below, text width=140mm, align=center] at (22.5,-0.1) {每级 9 位 = 512 项；L1 一项管 1\,GiB，L2 一项管 2\,MiB，L3 一项管 4\,KiB};
  \end{tikzpicture}
  \caption{39 位虚拟地址的划分}
  \label{fig:va}
\end{figure}

\file{kernel/aarch64.h} 用两个宏取出各级索引：

\begin{ccode}
#define PXMASK          0x1FF                    // 9 位
#define PXSHIFT(level)  (39-(level)*9)           // L1: 30, L2: 21, L3: 12
#define PX(level, va)   ((((uint64)(va)) >> PXSHIFT(level)) & PXMASK)
\end{ccode}

第 1、2 级表项可以是\emph{表描述符}（指向下一级页表）或\emph{块描述符}（直接映射 1\,GiB/2\,MiB 的连续物理内存）；第 3 级表项是\emph{页描述符}。启动页表用 2\,MiB 块，最终页表全部用 4\,KiB 页。

\subsection{硬件如何遍历页表}
\label{sec:hw-walk}

虚拟内存给每个进程一个错觉：整个地址空间都归它一个人用。进程发出的每个地址都是虚拟地址，由 MMU（内存管理单元）在访存的路上翻译成物理地址，进程本身完全察觉不到。这样做同时解决了两个问题：一是\emph{隔离}，进程只能访问自己页表里映射了的内存，碰不到别人的数据，也碰不到内核；二是\emph{分配简单}，每个进程的程序都可以从地址 0 开始链接，不必关心物理内存里哪些地方已经被别人占了。

MMU 翻译一个地址的过程叫\emph{页表遍历}（table walk），完全由硬件完成，软件只负责事先把页表在内存里摆好、再把根页表的物理地址写进 \reg{TTBR0\_EL1} 或 \reg{TTBR1\_EL1}。以用户地址为例：

\begin{enumerate}
  \item 从 \reg{TTBR0\_EL1} 取出 L1 页表的\textbf{物理}地址；
  \item 用 VA[38:30] 作下标读出 L1 表项。若它是表描述符，其中 [47:12] 位就是 L2 页表的物理地址；
  \item 用 VA[29:21] 在 L2 页表中读出表项，同理得到 L3 页表的物理地址；
  \item 用 VA[20:12] 在 L3 页表中读出页描述符，得到物理页框号；
  \item 物理页框号拼上 VA[11:0] 的页内偏移，就是最终的物理地址。
\end{enumerate}

每一步读到的地址都是物理地址——MMU 遍历页表时不可能再去翻译页表自己的地址。这就是为什么 \code{walk()} 往表项里写的是 \code{V2P(pagetable)}，读出来以后又要 \code{P2V()} 才能被内核软件访问（\ref{sec:walk-mappages} 节）。途中任何一级遇到无效表项（bit~0 为 0），遍历就停下来，CPU 产生一个\emph{翻译错误}（translation fault），并把出错的级别记在 \reg{ESR\_EL1} 里、出错的虚拟地址记在 \reg{FAR\_EL1} 里。

\paragraph{块映射.} 有时候要映射一大片连续的物理内存，例如启动时的整个内核镜像。这时可以让 L2 表项直接是一个\emph{块描述符}，指向一个 2\,MiB 对齐的物理块，遍历在第 3 步就结束：省去一级页表，页内偏移变成 VA[20:0] 共 21 位（$2^{21}=2$\,MiB）。同理，L1 块描述符映射 1\,GiB。描述符最低两位区分这几种情况：在 L1/L2 上 \code{01} 是块、\code{11} 是表；在 L3 上只有 \code{11} 合法，表示一个 4\,KiB 页——同样是 \code{11}，含义随所在的级别而变。xv6-aarch64 只在启动页表里用块映射（\ref{sec:bootpgt} 节）。

\paragraph{实例：一次真实的三级翻译.} 图~\ref{fig:mmu-walk} 用一个进程快照里的真实页表项，把上面五步走了一遍：pid 16 的用户栈地址 \code{0x89d0}。\code{TCR\_EL1.T0SZ = 25}，即 39 位 VA。VA[63:39] 全为 0，选 \reg{TTBR0\_EL1} = \code{0x6eb8000}；L1、L2、L3 的下标分别是 0、0、8，页内偏移是 \code{0x9d0}。MMU 依次读三个物理地址：
\begin{itemize}
  \item 读 \code{0x6eb8000}，得到表描述符 \code{0x6eb5003}；
  \item 读 \code{0x6eb5000}，得到表描述符 \code{0x6ec0003}；
  \item 读 \code{0x6ec0040}，得到页描述符 \code{0x6ece74b}。
\end{itemize}
页框 \code{0x6ece000} 拼上偏移，得到物理地址 \code{0x6ece9d0}。描述符低 12 位 \code{0x74b} 的含义见下一小节：AttrIndx=2 表示普通可缓存内存，AP=01 表示 EL0 可读写，SH=11，AF=1，UXN 和 PXN 都是 0。

\begin{figure}[htbp]
  \centering
  \includegraphics[width=\linewidth]{mmu-walk.pdf}
  \caption{MMU 三级页表翻译实例：VA \code{0x89d0} $\to$ PA \code{0x6ece9d0}（数值取自进程快照）}
  \label{fig:mmu-walk}
\end{figure}

\subsection{描述符的属性位}

表~\ref{tab:pte-bits} 是项目中用到的属性位。

\begin{table}[htbp]
  \centering
  \caption{页表项属性位（\file{kernel/aarch64.h}）}
  \label{tab:pte-bits}
  \small
  \begin{tabularx}{\linewidth}{@{}llX@{}}
    \toprule
    宏 & 位 & 含义 \\
    \midrule
    \code{PTE\_VALID} / \code{PTE\_TABLE} & [1:0] & L1/L2：\code{01} 块，\code{11} 表；L3：\code{11} 页（\code{PTE\_V}） \\
    \code{PTE\_INDX(i)} & [4:2] & AttrIndx：选择 \reg{MAIR\_EL1} 中的第 $i$ 个内存类型 \\
    \code{PTE\_U}、\code{PTE\_RO} & [7:6] & AP：EL0 是否可访问、是否只读 \\
    \code{PTE\_SH\_INNER} & [9:8] & 共享属性：\code{11} 为 Inner Shareable \\
    \code{PTE\_AF} & [10] & Access Flag：必须置 1，否则首次访问触发 Access Fault \\
    \code{PTE\_PXN}、\code{PTE\_UXN} & 53、54 & 内核态 / 用户态禁止执行 \\
    \bottomrule
  \end{tabularx}
\end{table}

\subsection{内存属性间接表}

页表项里并不直接写“可缓存”“设备内存”之类的属性，而是写一个 3 位的索引，真正的属性放在 \reg{MAIR\_EL1} 的 8 个字节里。项目定义了三种类型：

\begin{ccode}
#define AI_DEVICE_nGnRnE_IDX  0x0
#define AI_NORMAL_NC_IDX      0x1
#define AI_NORMAL_IDX         0x2
#define MT_DEVICE_nGnRnE  0x00    // 设备内存：不聚合、不重排、无提前写确认
#define MT_NORMAL_NC      0x44    // 普通内存，不可缓存
#define MT_NORMAL         0xff    // 普通内存，Inner/Outer Write-Back

#define PTE_DEVICE    PTE_INDX(AI_DEVICE_nGnRnE_IDX)
#define PTE_NORMAL_NC PTE_INDX(AI_NORMAL_NC_IDX)
#define PTE_NORMAL   (PTE_INDX(AI_NORMAL_IDX) | PTE_SH_INNER)
\end{ccode}

用途划分很简单：RAM、内核、用户页面和页表本身用 \code{PTE\_NORMAL}；UART、GPIO、SD 控制器等外设寄存器用 \code{PTE\_DEVICE}；摄像头 DMA 缓冲区用 \code{PTE\_NORMAL\_NC}，这样就不必在每帧前后做缓存维护（第~\ref{chap:drivers} 章）。

\section{启动页表}
\label{sec:bootpgt}

第~\ref{chap:boot} 章的启动代码把处理器降到 EL1、点亮串口之后，内核还运行在物理地址上，MMU 和缓存都是关着的。本节接着讲从这里到跳进高地址的三步：清 BSS、建临时页表、打开 MMU。

\subsection{清 BSS 与早期栈}

C 语言假定未初始化的全局变量为 0，这些变量位于 \code{.bss} 段。镜像文件里并不包含 BSS 的内容，所以要先清零。链接脚本导出了 \code{bss\_start} 和以 8 字节为单位的 \code{bss\_size}：

\begin{asmcode}
entry:
        adrp x1, bss_start          // PC 相对：MMU 未开时得到物理地址
        ldr w2, =bss_size
1:      cbz w2, 2f
        str xzr, [x1], #8
        sub w2, w2, #1
        b 1b
2:      ...                         // 输出 'B'
\end{asmcode}

把需要清零的数据单独放在 \code{.bss} 段，是为了节省镜像空间：ELF 只记录这个段的大小，不存它的内容。代价是装载之后必须由启动代码自己清零，所以链接脚本要导出段的起止位置，并让起点按 8 字节对齐。对齐在这里不只是为了效率：MMU 关闭时所有数据访问都按 Device 内存处理，非对齐访问会直接触发对齐异常，因此一次写 8 字节的 \code{str xzr} 要求地址必须是 8 的倍数。

栈也是同样的道理：C 代码需要一个可用的栈指针。最简单的做法是把 \code{sp} 设到内核镜像上方一个“足够远”的固定地址（例如内核装在 0 时设为 4\,MiB），栈向下增长也碰不到内核。本项目要支持四个核同时运行，所以在 BSS 中定义了 \code{stack0[4096 * NCPU]}，每个核按核号取其中一页（见 \ref{sec:mmu-on} 节末尾的 \code{\_start}）。清完 BSS 以后，主核先临时把 \code{sp} 设成 \code{stack0 + 4096}（这时还是物理地址），因为下一步建启动页表要调用一个 C 函数。

\subsection{两棵临时页表}


MMU 打开前要先建两棵只有两级的小页表，四张页表页（\code{l1entrypgt}、\code{l2entrypgt}、\code{l1kpgt}、\code{l2kpgt}）静态定义在 \file{start.c} 中、按页对齐。

\begin{figure}[htbp]
  \centering
  \includegraphics[width=\linewidth]{boot-el1-memory.png}
  \caption{启动页表：同一段物理内存同时映射到低地址和高地址。TTBR0\_EL1 恒等映射 VA $=$ PA，只在开 MMU 的过渡期使用；TTBR1\_EL1 映射 VA $=$ \code{KERNBASE} $+$ PA。图中按通用写法标的高地址 \code{0xffff0000\_\ldots}，在 xv6 里一律读作 \code{KERNBASE} $=$ \code{0xffffff80\_00000000} 加物理地址，例如物理 \code{0x80000} 对应 \code{0xffffff80\_00080000}，外设 \code{0x3f000000} 对应 \code{0xffffff80\_3f000000}}
  \label{fig:bootpgt}
\end{figure}

如图~\ref{fig:bootpgt}，恒等映射保证“打开 MMU 的那条指令之后还能取到下一条指令”，高地址映射保证“跳到链接地址之后一切照常”。两者都用 2\,MiB 块描述符，所以每棵树只要一张 L1 和一张 L2。早期版本在 \file{entry.S} 里用汇编循环填表；现在 \file{entry.S} 只把四张表的物理地址和 \code{EARLYTOP} 放进 \code{x0}～\code{x4}，调用 \file{start.c} 里的 C 函数 \code{entrypgtinit()}。MMU 没开时 C 代码同样运行在物理地址上，只要它不访问全局变量的绝对地址，就和汇编没有区别：

\begin{ccode}
void entrypgtinit(pte_t *l1entry, pte_t *l2entry, pte_t *l1kernel, pte_t *l2kernel,
                  uint64 pa_end)                    // pa_end = EARLYTOP
{
  const uint64 l2block = 1UL << PXSHIFT(2);         // 2 MiB
  for(uint64 pa = 0; pa < pa_end; pa += l2block)    // 恒等映射
    l2entry[PX(2, pa)] = PA2PTE(pa) | PTE_AF | PTE_NORMAL | PTE_VALID;
  l1entry[PX(1, 0)] = PA2PTE((uint64)l2entry) | PTE_TABLE | PTE_VALID;

  for(uint64 pa = 0; pa < pa_end; pa += l2block)    // KERNBASE + PA
    l2kernel[PX(2, KERNBASE + pa)] = PA2PTE(pa) | PTE_AF | PTE_NORMAL | PTE_VALID;
  uint64 uart_pa = 0x3f200000UL;                    // Mini UART 所在的 2 MiB
  l2kernel[PX(2, KERNBASE + uart_pa)] = PA2PTE(uart_pa) | PTE_AF | PTE_DEVICE | PTE_VALID;
  l1kernel[PX(1, KERNBASE)] = PA2PTE((uint64)l2kernel) | PTE_TABLE | PTE_VALID;
}
\end{ccode}

启动页表覆盖物理内存的前 4\,MiB（\code{EARLYTOP}$=$\code{0x400000}，两个 2\,MiB 块），外加 Mini UART 所在的一个设备块，供 \code{boot\_uart\_mark()} 在开 MMU 之后继续输出标记。\code{EARLYTOP} 原来是 2\,MiB，只盖住内核镜像；后来内核的 BSS 变大，\code{kinit1()} 在 \code{[end, 2\,MiB)} 里只剩 65 页，而 \code{kvminit()} 逐页映射 112\,MiB 内存需要约 70 个页表页，启动时就 \code{panic: kvmmap}，所以改成 4\,MiB。\code{main()} 一开始仍检查内核镜像连同 BSS 没有越过这个边界——否则后面的分配器会把内核自己的内存当成空闲页发出去：

\begin{ccode}
if(V2P(end) > EARLYTOP){
  for(char *m = "\nkernel image extends past EARLYTOP\n"; *m; m++)
    boot_uart_mark(*m);
  for(;;) asm volatile("wfe");
}
\end{ccode}

\subsection{打开 MMU：TCR、MAIR 与 SCTLR}
\label{sec:mmu-on}

页表摆好以后，还要告诉 MMU 怎样解释它们。主核在 \code{\_entry}、副核在 \code{entryothers} 里执行的是同一段代码：

\begin{asmcode}
        adrp x0, l1entrypgt
        adrp x1, l1kpgt
        msr  ttbr0_el1, x0        // 低地址：恒等映射
        msr  ttbr1_el1, x1        // 高地址：内核
        ldr  x0, =(TCR_T0SZ(25)|TCR_IRGN0(1)|TCR_ORGN0(1)|TCR_SH0(3)|TCR_TG0(0)| \
                   TCR_T1SZ(25)|TCR_IRGN1(1)|TCR_ORGN1(1)|TCR_SH1(3)|TCR_TG1(2)|TCR_IPS(0))
        msr  tcr_el1, x0
        ldr  x1, =((MT_DEVICE_nGnRnE<<(8*AI_DEVICE_nGnRnE_IDX)) | \
                   (MT_NORMAL_NC<<(8*AI_NORMAL_NC_IDX)) | (MT_NORMAL<<(8*AI_NORMAL_IDX)))
        msr  mair_el1, x1
        isb
        ldr  x1, =_start          // 链接地址（高地址）
        dsb  ish
        tlbi vmalle1              // 丢掉固件可能留下的翻译
        dsb  ish
        isb
        ic   iallu
        dsb  sy
        isb
        mrs  x0, sctlr_el1
        orr  x0, x0, #1           // M：打开 MMU
        orr  x0, x0, #(1 << 2)    // C：数据缓存
        orr  x0, x0, #(1 << 12)   // I：指令缓存
        msr  sctlr_el1, x0
        isb
        br   x1                   // 跳到高地址
\end{asmcode}

\reg{TCR\_EL1}（Translation Control Register）的各个字段决定了前面所有的“规则”：

\begin{itemize}
  \item \code{T0SZ}$=$\code{T1SZ}$=25$：两个地址空间都是 39 位，于是遍历从 L1 开始、一共三级；
  \item \code{TG0}$=0$、\code{TG1}$=2$：两边都是 4\,KiB 粒度。两个字段的编码不一样（TG0 中 \code{00} 是 4\,KiB，TG1 中 \code{10} 才是 4\,KiB），抄错一个数字，高地址的页表就会被按 64\,KiB 或 16\,KiB 粒度解释；
  \item \code{IRGN}/\code{ORGN}$=1$、\code{SH}$=3$：MMU 遍历页表时本身也要读内存，这几位规定这些读取走 Write-Back 缓存、属于 Inner Shareable 域。它们必须和页表页本身的映射属性（\code{PTE\_NORMAL}）一致，否则 CPU 写进缓存的新表项，MMU 可能看不到（第~\ref{sec:fault1} 节的故障与此同源）；
  \item \code{IPS}$=0$：物理地址 32 位（4\,GiB），对 BCM2837 足够。
\end{itemize}

\reg{MAIR\_EL1} 把三种内存类型写进第 0、1、2 字节，页表项的 AttrIndx 就是这里的下标。\reg{SCTLR\_EL1} 是总开关：M 位打开地址翻译，C、I 位打开数据和指令缓存。打开前的 \code{tlbi vmalle1} 和 \code{ic iallu} 清掉固件阶段可能残留的 TLB 和指令缓存；写 \reg{SCTLR\_EL1} 之后的 \code{isb} 保证后续取指都已经按新的翻译进行。此时 PC 仍是低地址，靠 \reg{TTBR0} 的恒等映射继续执行，直到 \code{br x1} 跳进高地址——这就是启动页表要同时建两棵树的原因。

\subsection{跳到高地址}

\code{ldr x1, =\_start} 从文字池中取出链接时确定的高虚拟地址；这里不能写 \code{b \_start}：打开 MMU 并不会改变 PC，而 \code{b} 和 \code{adr} 都是相对当前 PC 计算目标的，会落到“没开 MMU 时 \code{\_start} 所在的低地址”。必须事先用 \code{ldr} 从文字池取出链接器算好的高地址绝对值；这条 \code{ldr} 本身也是 PC 相对的，所以要在打开 MMU 之前执行。\code{br x1} 之后，CPU 就运行在 \reg{TTBR1\_EL1} 映射的高地址上了。\code{\_start} 为每个核在 \code{stack0} 中划出 4\,KiB 栈，然后跳进 C 函数 \code{main()}：

\begin{asmcode}
_start:
        ldr x0, =stack0
        mov x1, #1024*4
        mrs x2, mpidr_el1
        and x2, x2, #0x3
        add x2, x2, #1
        mul x1, x1, x2
        add x0, x0, x1               // sp = stack0 + (cpuid + 1) * 4096
        mov sp, x0
        mov w0, #'H'
        bl boot_uart_mark
        b main
\end{asmcode}

\section{多核启动}
\label{sec:smp-boot}

四个核在上电后都会执行 \code{\_entry}——CPU0 由 armstub 直接跳入，CPU1～3 则要等主核在 spin-table 里写入地址。如果像 \code{kernel\_old=1} 那样没有 armstub，四个核会同时从同一个入口开始执行，内核要立刻用 \reg{MPIDR\_EL1} 的低 8 位区分核号，把除 0 号核以外的所有核挂进死循环——这是单核教学内核的标准做法，但被挂起的核就再也没有机会醒来了。本项目需要四核同时调度，所以副核不能死循环，而要在一个可以被唤醒的地方等待。整个过程有两道“闸门”：

\paragraph{第一道闸门：启动页表。} 副核到达 \code{\_entry} 后，统一好异常级别（\ref{sec:xv6-el} 节）就进入 \code{secondary\_wait}，用 \code{ldar}（带 acquire 语义的加载）读取标志 \code{secondary\_release}，为 0 就 \code{wfe} 休眠。主核建好启动页表后，用 \code{stlr}（带 release 语义的存储）把标志置 1，再执行 \code{sev} 唤醒所有在 \code{wfe} 上休眠的核：

\begin{asmcode}
        adrp x10, secondary_release            // 主核
        add x10, x10, :lo12:secondary_release
        mov w11, #1
        stlr w11, [x10]
        sev

secondary_wait:                                // 副核
        adrp x10, secondary_release
        add x10, x10, :lo12:secondary_release
4:      ldar w11, [x10]
        cbnz w11, entryothers
        wfe
        b 4b
\end{asmcode}

acquire/release 配对保证：副核看到标志为 1 时，主核在此之前写入的页表内容也一定可见。

\paragraph{第二道闸门：内核初始化。} 主核在 \code{main()} 的最后调用 \code{start\_secondary\_cpus()}，把 \code{\_entry} 的物理地址写进三个 spin-table 槽，用 \code{dc cvac} 把这几行缓存写回内存（此时副核的 MMU 和缓存都没打开，只能看到内存里的值），再 \code{sev}：

\begin{ccode}
static void start_secondary_cpus(void) {
  volatile uint64 *spin = (volatile uint64*)P2V(0xe0);
  uint64 entry = V2P((uint64)_entry);
  for(int cpu = 1; cpu < NCPU; cpu++){
    spin[cpu - 1] = entry;
    asm volatile("dc cvac, %0" :: "r"(&spin[cpu - 1]) : "memory");
  }
  asm volatile("dsb sy\n\tsev" ::: "memory");
}
\end{ccode}

副核进入 \code{main()} 后在 \code{while(started == 0)} 处等待主核完成全部初始化，然后只做每核相关的部分：切换到最终页表、安装异常向量、打开本核的定时器中断路由，最后进入调度器。真机串口上会看到：

\begin{txtcode}
PH0hart 3 starting
hart 2 starting
hart 1 starting
\end{txtcode}

三个副核几乎同时出发，谁先抢到 \code{printf} 的锁谁先打印，所以顺序不固定；\code{PH0} 是不经过锁的早期标记，会与其他输出交错。

\section{从设备树到 memblock}
\label{sec:memblock}

建立最终页表需要分配页表页，分配页表页又需要物理页分配器，而物理页分配器的 \code{mem\_map} 又要按 RAM 的大小来开——这是一个“先有鸡还是先有蛋”的问题。早先的 xv6 用常量回避它：\code{PHYSTOP}$=$112\,MiB 是 RAM 的大小，\code{kinit1()} 先把 \code{[end, EARLYTOP)} 交给分配器建页表，\code{kinit2()} 再交出 \code{[EARLYTOP, PHYSTOP)}。树莓派 3 实际有 1\,GiB 内存（扣掉 GPU 部分约 960\,MiB），常量让八成内存闲置；换一块板子还得改源码重编。

Linux 的回答分两步：先从固件拿到内存布局（x86 的 e820、UEFI 内存表，ARM 上是设备树），记进 \code{memblock} 这个只有两张区间表的启动期分配器；等页表和 \code{struct page} 数组都就绪，\code{memblock\_free\_all()} 再把剩下的空闲内存一次性移交给伙伴系统。当前源码树照这个形状实现了一个迷你版本（\file{kernel/mm/fdt.c}、\file{kernel/mm/memblock.c}，合计约 500 行），\code{main()} 的开头变成：

\begin{ccode}
memblock_init(boot_dtb_pa);  // 从 DTB 读出 RAM 与保留区
kvminit();                   // 最终内核页表，页表页来自 memblock_alloc()
kvminithart();               // 切换 TTBR1，作废 TLB
page_alloc_init();           // mem_map 也来自 memblock，然后 memblock_free_all()
\end{ccode}

\subsection{设备树从哪里来}

AArch64 的启动约定是：引导程序把扁平设备树（flattened device tree，FDT，也叫 DTB）的物理地址放在 \code{x0} 里跳进内核。树莓派固件按 \file{config.txt} 加载 \file{bcm2710-rpi-3-b.dtb}，填好 \code{/memory} 节点后传给内核；QEMU 的 \code{-dtb} 选项做同样的事，并把 \code{/memory} 改写成模拟的内存大小。\file{entry.S} 的第一条指令就把 \code{x0} 存进 \code{x21}——后面清 BSS、建启动页表都不能碰它——等启动页表建好后再写进全局变量 \code{boot\_dtb\_pa}。

此时只有启动页表，它只覆盖物理前 4\,MiB，而 QEMU 把 DTB 放在 \code{0x08000000}，固件则通常放在 RAM 顶部附近。所以 \code{entrypgtinit()} 多收一个参数：DTB 地址非零、8 字节对齐、低于外设窗口时，把它所在的 2\,MiB 块和下一块（DTB 不超过 2\,MiB，可能跨块）也映射进高地址的 \code{l2kernel}，\code{memblock\_init()} 就能经 \code{P2V()} 直接读它。

\subsection{解析 FDT}

FDT 是一段大端格式的二进制：40 字节的头（魔数 \code{0xd00dfeed}、总长、结构块和字符串块的偏移、版本），之后是 \code{/memreserve/} 表（若干 \code{<地址, 长度>} 对，以 0 结尾）、结构块和字符串块。结构块是一串 32 位的 token：

\begin{ccode}
FDT_BEGIN_NODE 1   // 后跟以 NUL 结尾、4 字节对齐的节点名
FDT_END_NODE   2
FDT_PROP       3   // 后跟 len、nameoff（指向字符串块）和 len 字节的值
FDT_NOP        4
FDT_END        9
\end{ccode}

\code{fdt\_scan()} 不建树，只顺序走一遍 token，用一个深度不超过 8 的小栈记住每层的 \code{\#address-cells}/\code{\#size-cells}——它们描述的是\emph{子节点} \code{reg} 的格式，所以解析一个节点的 \code{reg} 要用父节点的值，而且属性可能出现在 \code{reg} 之后，于是 \code{reg} 留到 \code{FDT\_END\_NODE} 才上报。它关心三类节点：
\begin{itemize}
  \item 根下名为 \code{memory} 或 \code{memory@...}、或 \code{device\_type = "memory"} 的节点：\code{reg} 里每一对都回调 \code{memblock\_add()}；
  \item \code{/reserved-memory} 的子节点：有 \code{reg} 的回调 \code{memblock\_reserve()}；只有 \code{size} 的（如 CMA 池）是“请内核找个地方”，迷你版不支持，跳过；
  \item 根节点的 \code{model} 字符串，仅用于打印。
\end{itemize}
\code{/memreserve/} 表里的区间同样进 \code{reserved}（树莓派的 DTB 用它保护物理第 0 页的 armstub 与 spin-table）。头部校验失败（魔数不对、地址不对齐、长度超过 2\,MiB、版本低于 16）或扫描中遇到越界 token，\code{memblock\_init()} 就丢掉已经加进去的区间，退回内置布局。

\subsection{memblock：两张有序区间表}

\begin{ccode}
struct memblock_region { uint64 base, size; };
struct memblock_type   { int cnt; uint64 total; char *name;
                         struct memblock_region regions[32]; };
struct memblock { uint64 current_limit;
                  struct memblock_type memory, reserved; };
\end{ccode}

两张表都按地址排序，\code{memblock\_add()} 和 \code{memblock\_reserve()} 插入时与重叠或相邻的区间合并，\code{memblock\_remove()} 必要时把一个区间劈成两段。表是静态数组：此时还没有任何分配器，Linux 也一样先用静态的 128 项数组。\code{memblock\_init()} 在 DTB 之外再登记几样东西：

\begin{ccode}
if(memblock.memory.cnt == 0){           // 没有 DTB，或 DTB 不可用
  memblock_source = "built-in";
  memblock_add(0, PHYSTOP);             // 退回原来的 112 MiB
}
memblock_remove(PERIPHERAL_BASE_PA, ~0ULL - PERIPHERAL_BASE_PA);
memblock_reserve(0, PGROUNDUP(V2P(end)));                  // 第 0 页与内核镜像
memblock_reserve(CARVEOUT_START, CARVEOUT_END - CARVEOUT_START); // ramdisk/摄像头 DMA/模块区
memblock_set_current_limit(EARLYTOP);   // 现在只有前 4 MiB 有映射
\end{ccode}

\code{PHYSTOP} 因此从“RAM 的大小”降格为“没有 DTB 时的内置大小”。\code{0x07000000} 起的 6\,MiB 窗口原先在 \code{PHYSTOP} 之上，分配器天然碰不到；现在 RAM 远大于它，就得显式保留。DTB 本身也被保留，以后还可以再读。

\code{memblock\_alloc(size, align)} 和 Linux 一样\textbf{自顶向下}找：从最高的 memory 区间、\code{current\_limit} 以下开始，取对齐的候选地址；若与某个 reserved 区间重叠，就退到那个区间的下方再试。分到的区间立即记进 \code{reserved}，所以它永远不会再交给伙伴系统。\code{current\_limit} 是这里的关键：建最终页表时它等于 \code{EARLYTOP}，保证页表页落在启动页表已经映射的 4\,MiB 里；切换到最终页表后它抬到 RAM 顶部，\code{mem\_map} 就可以放到任何地方。

\section{物理内存初始化与伙伴系统}
\label{sec:buddy}

在伙伴系统就绪之前，\code{kalloc()} 直接转给 \code{memblock\_alloc(PGSIZE, PGSIZE)}，\code{kvmmake()} 和 \code{walk()} 不用知道背后是谁。内核页表永不释放，这些页留在 \code{reserved} 里正合适；在 \code{buddy\_ready} 置位之前调用 \code{kfree()} 会直接 \code{panic}。

\subsection{从空闲链表到伙伴系统}

最初的分配器就是经典 xv6 的空闲链表：每个空闲页的前 8 字节存放下一页的地址，\code{kalloc()} 从表头取，\code{kfree()} 往表头放。它简单可靠，但有三个缺陷：

\begin{itemize}
  \item \textbf{只能分配单页}。链表里相邻的两项在物理上毫无关系，拿不到连续的多页。于是需要连续物理内存的地方只能另想办法：virtio 驱动用静态数组 \code{pages[2*PGSIZE]}，\code{init\_module} 用静态的 \code{module\_image[]}，摄像头 DMA 缓冲和模块区则干脆划在当时的 RAM 上限 \code{PHYSTOP} 之上，绕开分配器。
  \item \textbf{页没有元数据}。分配器不知道一页是空闲还是在用、被谁用，重复 \code{kfree()} 同一页不会报错，只会让链表成环，很久以后才以莫名其妙的方式崩溃。也没有地方放引用计数，写时复制无从谈起。
  \item \textbf{统计要遍历}。\code{/proc/meminfo} 的 MemFree 靠 \code{kfreepages()} 把整条链表走一遍。
\end{itemize}

现在的 \file{kernel/mm/kalloc.c} 换成了 Linux 式的两件套：为每个物理页框准备一个描述符 \code{struct page}，组成数组 \code{mem\_map}；空闲内存按 $2^{order}$ 页的块组织在 \code{free\_area[order]} 里，也就是伙伴系统。\code{kalloc()}/\code{kfree()} 的接口原样保留，变成 order 0 的薄封装，全内核四十多处调用点一行都不用改。

\subsection{\code{struct page} 与 \code{mem\_map}}

\file{include/kernel/page.h} 定义了描述符：

\begin{ccode}
struct list_head {
  struct list_head *next, *prev;
};

struct page {
  struct list_head lru;         // PG_buddy：挂在 free_area[order] 上
                                // PG_pcp：挂在每 CPU 链表上
                                // PG_slab：挂在 kmem_cache 的 partial 链表上
  uint8  flags;                 // PG_*
  uint8  order;                 // 空闲块或多页分配的阶
  uint16 inuse;                 // PG_slab：这个 slab 已分出的对象数
  int    refcount;              // 0 空闲，分配后为 1
  struct kmem_cache *slab_cache;// PG_slab：所属的 cache
  void  *freelist;              // PG_slab：这个 slab 的第一个空闲对象
};

#define PG_reserved 0x01        // 从不交给伙伴系统（内核镜像、mem_map 本身）
#define PG_buddy    0x02        // 空闲块的首页，挂在 free_area[] 上
#define PG_head     0x04        // order > 0 分配块的首页
#define PG_slab     0x08        // 属于某个 kmem_cache（见 slab 一节）
#define PG_pcp      0x10        // 空闲，暂存在每 CPU 链表上
\end{ccode}

一个描述符 40 字节（最初只有前 24 字节，第~\ref{sec:slab} 节为 slab 加了后两个字段；Linux 是 64 字节，同样按页的用途复用字段），RAM 里每 4\,KiB 一个，物理页号 \code{pfn} 就是数组下标。地址换算因此都是常数时间：

\begin{ccode}
struct page *pfn_to_page(uint64 pfn) { return mem_map + pfn; }
uint64 page_to_pfn(struct page *pg)  { return pg - mem_map; }
struct page *virt_to_page(void *kva) { return pfn_to_page(V2P(kva) >> PGSHIFT); }
void *page_address(struct page *pg)  { return P2V(page_to_pfn(pg) << PGSHIFT); }
\end{ccode}

\code{mem\_map} 自己也要占物理内存，而它的长度由 \code{memblock\_end\_of\_DRAM()} 决定：QEMU 给出 960\,MiB 时是 $245\,760$ 个描述符、9\,600\,KiB，没有 DTB 时是 $28\,672$ 个、1\,120\,KiB。9.4\,MiB 已经放不进启动页表的 4\,MiB，所以 \code{page\_alloc\_init()} 放在 \code{kvminithart()} 之后——此时最终页表已经线性映射了全部 RAM，可以先抬高 \code{current\_limit}，再像 Linux 的 \code{alloc\_node\_mem\_map()} 那样从 memblock 里要一块：

\begin{ccode}
max_pfn = memblock_end_of_DRAM() >> PGSHIFT;
size = max_pfn * sizeof(struct page);
memblock_set_current_limit(max_pfn << PGSHIFT);  // 最终页表已覆盖全部 RAM
if((pa = memblock_alloc(PGROUNDUP(size), PGSIZE)) == 0)
  boot_halt("\npage_alloc_init: no memory for mem_map\n");
mem_map = (struct page *)P2V(pa);

memset(mem_map, 0, size);
for(uint64 pfn = 0; pfn < max_pfn; pfn++)
  mem_map[pfn].flags = PG_reserved;      // 先全部保留
memblock_free_all(free_range);           // 只有空闲区间才清掉 PG_reserved
buddy_ready = 1;
\end{ccode}

自顶向下的分配把 \code{mem\_map} 放在 RAM 的最高处（QEMU 下是 \code{0x3b6a0000}），低端留给 DMA 等可能在意地址的用途，这也是 Linux 的习惯。

\subsection{伙伴分配与合并}

空闲内存按块组织。\code{free\_area[order]} 是一条双向链表，挂着所有大小为 $2^{order}$ 页的空闲块，\code{MAX\_ORDER}$=11$，块大小从 4\,KiB 到 4\,MiB。每个块都\textbf{自然对齐}：首页的 pfn 是 $2^{order}$ 的倍数。只有块的首页带 \code{PG\_buddy} 并挂在链表上，块内其余页的描述符不参与。

自然对齐带来了伙伴系统名字的由来：一个 order 阶的块，与它“结伴”的另一半的首页号只差第 order 位，即

\[
  \text{buddy\_pfn} = \text{pfn} \oplus 2^{order}.
\]

\textbf{分配}时从请求的阶往上找第一个非空链表，取下一块；如果块比需要的大，就不断对半切开，把上半块挂回低一阶的链表（Linux 里叫 \code{expand()}）：

\begin{ccode}
for(o = order; o < MAX_ORDER; o++){
  if(list_empty(&zone.free_area[o].free_list))
    continue;
  pg = lru_to_page(zone.free_area[o].free_list.next);
  del_from_free_area(pg, o);
  while(o > order){             // 把用不到的上半块还回低阶链表
    o--;
    add_to_free_area(pg + (1UL << o), o);
  }
  break;
}
\end{ccode}

\textbf{释放}时反过来：只要伙伴也是同阶的完整空闲块，就把它从链表摘下、合成高一阶的块，一直合到伙伴不空闲或到达最高阶为止（\code{\_\_free\_one\_page()}）：

\begin{ccode}
while(order < MAX_ORDER - 1){
  uint64 buddy_pfn = pfn ^ (1UL << order);
  struct page *buddy;

  if(buddy_pfn >= max_pfn)
    break;
  buddy = pfn_to_page(buddy_pfn);
  if((buddy->flags & PG_buddy) == 0 || buddy->order != order)
    break;
  del_from_free_area(buddy, order);
  pfn &= ~(1UL << order);       // 合并后的块从低半块开始
  order++;
}
add_to_free_area(pfn_to_page(pfn), order);
\end{ccode}

\code{PG\_reserved} 页永远不带 \code{PG\_buddy}，合并自然在内核镜像和 \code{mem\_map} 的边界处停下，不需要任何特判。

图~\ref{fig:buddy-split} 以从一个空的 order 3 块里分配一页为例：8 页的块被切三次，留下 order 2、1、0 各一块；这一页释放时，三次合并按相反顺序发生，恢复成原来的 order 3 块。

\begin{figure}[htbp]
  \centering
  \begin{tikzpicture}[
      pg/.style={draw, minimum width=11mm, minimum height=6mm, inner sep=0pt, font=\scriptsize},
      lab/.style={font=\footnotesize, anchor=east},
      br/.style={decorate, decoration={brace, mirror, amplitude=3pt}}]
    \def\w{1.1}
    \foreach \i in {0,...,7}
      \node[pg, fill=cyan!15] at (\i*\w, 0) {\i};
    \node[lab] at (-0.8, 0) {order 3 空闲块};
    \node[pg, fill=orange!40] at (0, -1.2) {0};
    \node[pg, fill=green!20] at (\w, -1.2) {1};
    \foreach \i in {2,3}
      \node[pg, fill=yellow!30] at (\i*\w, -1.2) {\i};
    \foreach \i in {4,...,7}
      \node[pg, fill=cyan!15] at (\i*\w, -1.2) {\i};
    \node[lab] at (-0.8, -1.2) {alloc\_pages(0) 之后};
    \draw[br] (-0.5*\w,-1.65) -- (0.5*\w,-1.65) node[midway, below=3pt, font=\scriptsize] {已分配};
    \draw[br] (0.5*\w,-1.65) -- (1.5*\w,-1.65) node[midway, below=3pt, font=\scriptsize] {order 0};
    \draw[br] (1.5*\w,-1.65) -- (3.5*\w,-1.65) node[midway, below=3pt, font=\scriptsize] {order 1};
    \draw[br] (3.5*\w,-1.65) -- (7.5*\w,-1.65) node[midway, below=3pt, font=\scriptsize] {order 2};
    \node[font=\scriptsize, align=left, anchor=west] at (7.7*\w, -0.6)
      {伙伴：\\ $0\oplus1=1$\\ $0\oplus2=2$\\ $0\oplus4=4$};
  \end{tikzpicture}
  \caption{从 order 3 块分配一页：切三次，上半块依次挂回 order 2、1、0；释放时按 $\text{pfn}\oplus2^{order}$ 找伙伴逐级合并}
  \label{fig:buddy-split}
\end{figure}

\subsection{把内存交给伙伴系统}

\code{memblock\_free\_all()} 遍历 memory 减去 reserved 后剩下的每一段 \code{[start, end)}，对每段调用 \code{free\_range()}。最早的实现逐页 \code{kfree()}，每页还要先填一遍垃圾；新实现像 Linux 的 \code{\_\_free\_pages\_memory()} 那样，每次取当前位置能放下的\textbf{最大对齐块}直接挂上去：

\begin{ccode}
while(pfn < end_pfn){
  int order = MAX_ORDER - 1;
  while(order > 0 &&
        ((pfn & ((1UL << order) - 1)) != 0 || pfn + (1UL << order) > end_pfn))
    order--;
  for(uint64 i = 0; i < (1UL << order); i++){
    struct page *pg = pfn_to_page(pfn + i);
    pg->flags &= ~PG_reserved;
    pg->refcount = 0;
  }
  __free_one_page(pfn, order);  // 可能与上一段交出的块合并
  zone.managed_pages += 1UL << order;
  pfn += 1UL << order;
}
\end{ccode}

保留区把空闲内存切成几段，相邻两段交出的块在 \code{\_\_free\_one\_page()} 里照常合并。在 QEMU 上，受管理的页数是 \code{managed\_pages}$=241\,333$（约 943\,MiB），\code{/proc/meminfo} 的 MemTotal 报告的就是这个数，而不是 RAM 的总大小——与 Linux 的 MemTotal 只计 managed pages 一致。启动时 \code{memblock\_dump()} 打印两张表：

\begin{txtcode}
memblock: memory map from DTB at pa 0x0000000008000000 (89494 bytes),
          model "Raspberry Pi 3 Model B"
  memory.cnt = 1, total = 983040 KiB
    [0] 0x0000000000000000-0x000000003bffffff 983040 KiB
  reserved.cnt = 5, total = 17707 KiB
    [0] 0x0000000000000000-0x00000000001c3fff 1808 KiB
    [1] 0x00000000003ef000-0x00000000003fffff 68 KiB
    [2] 0x0000000007000000-0x00000000075fffff 6144 KiB
    [3] 0x0000000008000000-0x0000000008015d95 87 KiB
    [4] 0x000000003b6a0000-0x000000003bffffff 9600 KiB
buddy: 241333 pages managed, 241333 free,
       mem_map 9600 KiB at pa 0x000000003b6a0000
\end{txtcode}

五个保留区依次是：第 0 页加内核镜像与 BSS；最终内核页表的 17 个页表页（自顶向下，紧贴 \code{EARLYTOP} 之下）；ramdisk、摄像头 DMA 与模块区窗口；DTB 本身；\code{mem\_map}。不带 DTB 启动（\code{make qemu QEMU\_DTB=}）时第一行变成 \code{memory map from built-in}，RAM 回到 112\,MiB，受管理的约 2.8 万页。

\subsection{接口与防错}

对外的接口与 Linux 同名：

\begin{ccode}
struct page *alloc_pages(int order);          // 2^order 个物理连续页
void         __free_pages(struct page *pg, int order);
void        *__get_free_pages(int order);     // 返回内核虚拟地址
void         free_pages(void *kva, int order);

void *kalloc(void)   { return __get_free_pages(0); }
void  kfree(void *va);   // kalloc() 的页，或 kmalloc() 的内存
\end{ccode}

\code{kfree()} 把页填成 \code{0x01}、\code{kalloc()} 把页填成 \code{0x05} 的做法现在由编译开关 \code{PAGE\_POISONING} 控制（对应 Linux 的 \code{CONFIG\_PAGE\_POISONING}），slab 对象释放时填 \code{0x01} 也归它管。它默认关闭：每分配、释放一页都要把 4\,KiB 整个写一遍，fork、缺页、管道这些热路径都为此付出代价，而需要清零的地方（用户页、页表页、\code{kmem\_cache\_alloc()} 返回的对象）本来就各自显式 \code{memset}。追查内存错误时用 \code{make PAGE\_POISONING=1} 重新编译——第~\ref{chap:user} 章会看到，用户态异常地址 \code{0x0505050505050505} 正是一页刚分配、尚未初始化的内存，这类线索只有打开开关才会出现。有了描述符，释放路径还能做空闲链表做不到的检查，出错时直接打印物理地址并 \code{panic}：

\begin{itemize}
  \item 释放 \code{PG\_reserved} 页（内核镜像、\code{mem\_map}）；
  \item 重复释放：页已经带 \code{PG\_buddy} 或 \code{PG\_pcp}；
  \item 把 slab 页（\code{PG\_slab}）当普通页释放；
  \item 阶不匹配：用 order 3 分配、却按 order 0 释放，或反之（靠 \code{PG\_head} 和 \code{order} 判断）；
  \item 释放时 \code{refcount} 仍大于 1。
\end{itemize}

\code{refcount} 由 \code{get\_page()}/\code{put\_page()}/\code{page\_count()} 维护：分配时置 1，\code{put\_page()} 减到 0 时把页还给伙伴系统。这正是写时复制所需的那一份“这一页还有谁在用”，第~\ref{sec:cow-fork} 节的 COW fork 就建立在它之上。计数的加减用 \code{ldaxr}/\code{stlxr} 写成的 LL/SC 循环，而不是 GCC 的 \code{\_\_atomic} 内建函数：新版 aarch64 工具链默认把它们编译成对 \code{\_\_aarch64\_ldadd4\_acq\_rel} 等“外置原子操作”辅助函数的调用，而内核不链接 libgcc，链接时会报未定义符号。

\subsection{每 CPU 页缓存}
\label{sec:pcp}

四个核共用一把 \code{zone.lock}，而 \code{kalloc()}/\code{kfree()} 是内核里最频繁的调用之一：缺页、写时复制、管道和文件缓冲都要一页一页地拿和还。Linux 的办法是给每个 CPU 一个小的单页缓存（\code{struct per\_cpu\_pages}，简称 pcp）：大多数单页分配只碰本 CPU 的链表，\code{zone.lock} 每批才拿一次。xv6 照搬了这个形状：

\begin{ccode}
#define PCP_HIGH  64            // 超过这么多页就还一批给伙伴系统
#define PCP_BATCH 16            // 每次与伙伴系统之间搬运的页数

struct per_cpu_pages {
  struct spinlock lock;
  int count;
  struct list_head list;        // 空闲的 order-0 页，表头是最近释放的（最“热”）
} __attribute__((aligned(64))); // 各占一条 cache line，避免伪共享

static struct per_cpu_pages pcp[NCPU];
\end{ccode}

\begin{itemize}
  \item \textbf{分配}（\code{rmqueue\_pcplist()}）：本 CPU 的链表为空时，拿一次 \code{zone.lock}，用 \code{\_\_rmqueue(0)} 搬 \code{PCP\_BATCH} 页过来；然后从表头取一页。order $>0$ 的请求照旧直接走伙伴系统。
  \item \textbf{释放}（\code{free\_pages\_ok()} 的 order-0 分支，对应 Linux 的 \code{free\_unref\_page()}）：页挂到本 CPU 表头，标上 \code{PG\_pcp}；数量到 \code{PCP\_HIGH} 时，从表尾（最“冷”的一端）还 \code{PCP\_BATCH} 页给伙伴系统，让它们有机会合并。
  \item \textbf{排空}（\code{drain\_all\_pages()}）：任何分配失败时先把所有 CPU 的链表还给伙伴系统再试一次——别的核链表里的页，以及因此无法合并出的大块，都可能正好满足这次请求。
\end{itemize}

早期 Linux 靠“只有本 CPU 访问、访问时关中断”保护 pcp，排空别的 CPU 的链表得发 IPI 让它自己来做；5.19 以后每个 pcp 带一把自旋锁，平时只有本 CPU 拿，几乎不会争抢，而排空可以直接跨核进行。xv6 采用后一种：\code{this\_pcp\_lock()} 在关中断的情况下读 \code{cpuid()} 并拿本 CPU 的锁，锁的顺序固定为先 pcp、后 \code{zone.lock}。

\code{PG\_pcp} 让 pcp 上的页和伙伴系统里的页区分开：\code{\_\_free\_one\_page()} 只与带 \code{PG\_buddy} 的伙伴合并，不会把 pcp 链表上的页摘走；对 pcp 上的页再释放一次，同样会被当作重复释放抓住。

空闲页的统计也换成了计数器。原来的 \code{kfreepages()} 每次都拿着 \code{zone.lock} 把 11 个阶的 \code{nr\_free} 加一遍；现在 \code{add\_to\_free\_area()}/\code{del\_from\_free\_area()} 顺手维护 \code{zone.nr\_free}（对应 Linux 的 \code{NR\_FREE\_PAGES}），\code{kfreepages()} 不拿锁，返回它再加上各 CPU 的 \code{count}。pcp 上的页也算空闲，所以 \code{/proc/meminfo} 的 MemFree 不会因为页暂存在 pcp 上而“丢失”，cowtest 用 MemFree 做的泄漏检查依然精确；另有一行 PerCPUPages 显示 pcp 上的总量（四个核最多 $4\times 64$ 页，即 1\,MiB）。

\subsection{观察：启动自检与 \code{/proc/buddyinfo}}

\code{main()} 在 \code{printf} 可用之后立刻调用 \code{buddy\_selftest()}：先排空 pcp，记下各阶的空闲块数，分配 order 3、order 0 和 order 10 各一块，检查对齐、计数和引用计数，再全部释放并再次排空 pcp，要求每一阶的空闲块数都恢复原样——这同时验证了切分与合并。随后 \code{slab\_selftest()} 检查第~\ref{sec:slab} 节的 kmalloc。两者通过后各打印一行：

\begin{txtcode}
buddy: 241333 pages managed, 241333 free, mem_map 9600 KiB at pa 0x000000003b6a0000
slab: 7 kmalloc caches (32..2048 B), self-test ok, 1 page(s) kept by kmalloc-64
\end{txtcode}

\code{/proc/buddyinfo} 按 Linux 的格式列出 order 0 到 10 各有多少空闲块。下面是 112\,MiB 时的一次观察（带 DTB 启动时 order 10 一列是 230 多个）。刚登录时，几乎所有内存都还是 26 个完整的 4\,MiB 块；跑完 \code{mmaptest} 后，大块被切开，低阶链表上出现了大量碎块，order 10 的块少了一个：

\begin{shcode}
$ cat /proc/buddyinfo
Node 0, zone   Normal       1      0      0      0      1      1      1      1      1      1     26
$ mmaptest
...
mmaptest: OK (0 failures)
$ cat /proc/buddyinfo
Node 0, zone   Normal       0     10     18     18     17     17      9      3      0      0     25
\end{shcode}

表~\ref{tab:buddy-vs-linux} 把这一实现与 Linux 对照。

\begin{table}[htbp]
  \centering
  \caption{xv6 与 Linux 的物理页管理}
  \label{tab:buddy-vs-linux}
  \small
  \begin{tabularx}{\linewidth}{@{}lXX@{}}
    \toprule
    & Linux & xv6 \\
    \midrule
    内存探测 & e820 / UEFI / DTB & DTB 的 \code{/memory}；没有时用常量 \code{PHYSTOP} \\
    早期分配器 & memblock（memory / reserved 两张表） & 迷你 memblock，各 32 项静态数组 \\
    页描述符 & \code{struct page}，约 64 字节 & \code{struct page}，40 字节 \\
    \code{mem\_map} 来源 & \code{memblock\_alloc()} & \code{memblock\_alloc()}，自顶向下 \\
    节点与 zone & 多 NUMA 节点，DMA/DMA32/Normal & 单节点、单 zone \\
    阶数 & \code{MAX\_ORDER}（4\,KiB～4\,MiB） & 同为 11 阶 \\
    每 CPU 缓存 & \code{per\_cpu\_pages}，带锁（5.19 起） & 同样形状，high 64 / batch 16 \\
    空闲页统计 & \code{NR\_FREE\_PAGES} 计数器 & \code{zone.nr\_free} 加各 pcp 的 \code{count} \\
    小对象 & SLUB，\code{kmalloc-8} 到 \code{kmalloc-8k} & 简化的 SLUB，\code{kmalloc-32} 到 \code{kmalloc-2048} \\
    \bottomrule
  \end{tabularx}
\end{table}

物理 \code{0x07000000} 起的 6\,MiB 仍保留给 QEMU 时代的 ramdisk、摄像头 DMA 区和模块区，它们登记在 memblock 的保留表里，分配器永远不会把它们发出去。有了 \code{alloc\_pages(order)}，这类“必须物理连续”的缓冲以后可以改从伙伴系统里分配，或者像 Linux 那样写进设备树的保留内存节点，而不必在源码里手工划地盘。

\section{slab 与 \code{kmalloc()}}
\label{sec:slab}

伙伴系统的最小单位是一页。内核里大量对象只有几十到几百字节：一个 VMA 88 字节，一个 \code{struct file} 168 字节，伪终端的打开状态只有一个 \code{int}。为它们各要一页，浪费的是 95\% 以上的内存。xv6 原先已经有 \code{kmem\_cache}：每个 cache 从 \code{kalloc()} 拿页，切成等长的对象，用对象自身的第一个字串成空闲链表。但它有两个缺陷：

\begin{itemize}
  \item 只有\emph{类型专用}的 cache，没有按大小分配的 \code{kmalloc()}。于是 \file{unicam.c} 为 4 字节的 \code{struct cam\_file}、\file{pty.c} 为一个 \code{struct pty} 都直接 \code{kalloc()} 整页。
  \item 页一旦进了 cache 就永不归还。所有对象被释放后，空闲链表把不同页里的对象混在一起，cache 既不知道一个对象属于哪一页，也不知道某一页是否已经全空。
\end{itemize}

\subsection{元数据放进 \code{struct page}}

第二个缺陷的根源是空闲链表挂在 cache 上，而不是挂在页上。Linux 的 SLUB 把每个 slab（一页或几页）的元数据放在它的 \code{struct page} 里：哪个 cache、页内的空闲链表、已分出几个对象。xv6 照做，于是 \code{struct page} 多了 \code{slab\_cache}、\code{freelist}、\code{inuse} 三个字段（第~\ref{sec:buddy} 节的定义），\code{flags} 里多了 \code{PG\_slab}，\code{lru} 在 slab 页上用来挂 cache 的 partial 链表：

\begin{ccode}
struct kmem_cache {
  struct spinlock lock;
  const char *name;
  uint object_size, stride, alignment;
  uint objs_per_slab;           // PGSIZE / stride
  struct list_head partial;     // 还有空闲对象的 slab
  uint nr_partial;
  uint pages, total, inuse;     // 统计，供 /proc/slabinfo
  struct kmem_cache *next;      // 所有 cache 串成一条链
};
\end{ccode}

\begin{itemize}
  \item \textbf{分配}：取 partial 链表上第一个 slab，从它的 \code{freelist} 摘一个对象，\code{inuse++}；摘空了就把这个 slab 移出 partial（满的 slab 不挂在任何链表上，和 SLUB 一样）。partial 为空时先\emph{放开 cache 锁}再 \code{alloc\_pages(0)}——页分配器可能排空 pcp，调用者也可能还拿着别的锁——拿到页后在 \code{struct page} 里建好空闲链表，再挂进 partial。对象依旧在返回前清零。
  \item \textbf{释放}：\code{virt\_to\_page(obj)} 找到 slab，核对 \code{PG\_slab}、\code{slab\_cache} 和对象边界，然后把对象推回这个 slab 的 \code{freelist}。原来是满的就重新挂回 partial。\code{inuse} 减到 0 且 cache 还有别的 partial slab 时，清掉 \code{PG\_slab}，把页还给伙伴系统；只剩这一个时留着它，避免一对 alloc/free 让同一页在 slab 和伙伴系统之间来回搬（Linux 用 \code{min\_partial} 做同样的事）。
\end{itemize}

整页都用来放对象，没有页内头部，所以 \code{kmalloc-2048} 一页正好两个对象。

\subsection{尺寸类与 \code{kfree()}}

\code{kmalloc\_init()} 建 7 个通用 cache：\code{kmalloc-32}、\code{-64}、……、\code{-2048}，每个按自身大小对齐。\code{kmalloc(size)} 向上取到 2 的幂选 cache；超过 2\,KiB 的请求像 Linux 的 \code{kmalloc\_large()} 那样直接向伙伴系统要 $2^{order}$ 页。

释放只有一个入口 \code{kfree()}，它和 Linux 的 \code{kfree()} 一样靠页描述符判断内存从哪里来：

\begin{ccode}
void kfree(void *va) {
  struct page *pg = virt_to_page(va);
  if(pg->flags & PG_slab){            // kmalloc() 的小对象
    kmem_cache_free(pg->slab_cache, va);
    return;
  }
  if(pg->flags & PG_head){            // kmalloc() 的大块
    __free_pages(pg, pg->order);
    return;
  }
  free_pages(va, 0);                  // kalloc() 的一页
}
\end{ccode}

这样原有上百处 \code{kfree()} 调用一处不用改，\code{unicam.c}、\code{pty.c} 把 \code{kalloc()} 换成 \code{kmalloc(sizeof(*p))} 即可，释放端照旧 \code{kfree()}。\code{kmalloc} 也加进了模块的导出符号表，\file{.ko} 驱动可以直接用。

\subsection{观察：\code{/proc/slabinfo}}

每个 cache 在 \code{kmem\_cache\_init()} 时串进一条全局链表，\code{/proc/slabinfo} 按 Linux slabinfo 2.1 的列序打印它们。登录后不久的样子（节选）：

\begin{txtcode}
# name            <active_objs> <num_objs> <objsize> <objperslab> <pagesperslab> <num_slabs>
file_cache                    2         23       168           23              1           1
vma_cache                     3         42        88           42              1           1
proc_cache                   16         20       384           10              1           2
kmalloc-2048                  0          0      2048            2              1           0
kmalloc-64                    0          0        64           64              1           0
kmalloc-32                    0          0        32          128              1           0
\end{txtcode}

\code{num\_slabs} 为 0 的行说明 cache 此刻一页都没占：空 slab 已经还给了伙伴系统。启动时的 \code{slab\_selftest()} 分配 300 个 40 字节对象（落在 \code{kmalloc-64}，占 5 页），检查对齐、清零和页描述符，再全部释放，要求除 cache 自愿保留的一页外，空闲页数精确回到测试前的值。

\section{调试视图：\code{/proc/meminfo} 与 \code{/proc/iomem}}
\label{sec:proc-mem}

物理内存现在由 DTB、memblock、伙伴系统、pcp 和 slab 五层共同管理，出了问题首先要知道“内存都去哪儿了”。两个 procfs 文件按 Linux 的格式回答这个问题。

\code{/proc/meminfo} 给出各层的数量。MemTotal 是交给伙伴系统的页（managed pages），MemFree 是伙伴系统加上各 CPU pcp 里的空闲页，Slab 是所有 \code{kmem\_cache} 占用的页，Reserved 是 memblock 保留区与 RAM 的交集；DirectMap4k/DirectMap2M 遍历内核页表，统计线性映射里分别有多少字节是 4\,KiB 页、多少是 2\,MiB 块，可以直接验证第~\ref{sec:memblock} 节的块映射是否生效：

\begin{txtcode}
MemTotal:       965332 kB
MemFree:        963916 kB
MemAvailable:   963916 kB
Slab:           40 kB
PerCPUPages:    192 kB
Reserved:       17707 kB
DirectMap4k:    7684 kB
DirectMap2M:    972800 kB
CowShared:      76 pages
...
MemorySource:   DTB
\end{txtcode}

\code{/proc/iomem} 给出物理地址空间的布局：先按 memblock 的 memory 表列出 System RAM，再把 reserved 表里落在其中的区间缩进一层列在下面，并认出各自是什么——第 0 页到内核装载地址之间的 spin-table/armstub、内核代码与数据、启动时分配的内核页表、carve-out 窗口里的三块缓冲、DTB 本身、\code{mem\_map}。RAM 之后到外设窗口之间，带 DTB 时是 VideoCore 的那部分内存（帧缓冲 \code{0x3c100000} 就在这里），不带 DTB 时则如实标成“未描述”。最后是各外设控制器的 MMIO 子区间：

\begin{txtcode}
0x0000000000000000-0x000000003bffffff : System RAM
  0x0000000000000000-0x000000000007ffff : Reserved (spin-table, armstub)
  0x0000000000080000-0x00000000000bbfff : Kernel code
  0x00000000000bc000-0x00000000001c316f : Kernel data/bss
  0x00000000003ef000-0x00000000003fffff : Kernel page tables
  0x0000000007000000-0x00000000075fffff : Reserved (carve-out window)
    0x0000000007000000-0x00000000070fffff : QEMU xv6 ramdisk
    0x0000000007200000-0x00000000073fffff : camera DMA buffer (non-cacheable)
    0x0000000007400000-0x00000000074fffff : loadable-module arena (RWX)
  0x0000000008000000-0x0000000008015d95 : Flattened Device Tree
  0x000000003b6a0000-0x000000003bffffff : mem_map (struct page[])
0x000000003c000000-0x000000003effffff : Reserved (VideoCore / GPU memory, ...)
0x000000003f000000-0x000000003fffffff : BCM2837 peripheral MMIO window
  0x000000003f003000-0x000000003f003fff : BCM system timer
  ...
0x0000000040000000-0x0000000040000fff : BCM2836/7 ARM-local interrupt/timer MMIO
\end{txtcode}

两者都只读现成的数据结构（memblock 的两张表在启动后保留下来，正像 Linux 保留 \code{memblock} 供 \code{/sys/kernel/debug/memblock} 查看），不额外维护任何状态。

\subsection{逐条跟踪：\code{MM\_TRACE}}

\code{/proc} 给出的是快照；要看一次 \code{fork} 或一次 \code{sbrk} 之后内存到底是怎么一页页分出去的，需要逐事件的跟踪。与 \code{FS\_TRACE}、\code{VM\_TRACE} 一样，它是一个编译开关，定义在 \file{include/kernel/mm\_trace.h}：

\begin{shcode}
make MM_TRACE=1 qemu                      # 只跟踪用户进程（默认）
make MM_TRACE=1 MM_TRACE_PID=0 qemu       # 全部，包括内核线程和中断
make MM_TRACE=1 MM_TRACE_PID=7 MM_TRACE_LIMIT=500 qemu
\end{shcode}

关闭时 \code{MMTRACE(...)} 展开成 \code{if(0) printf(...)}：参数照样做类型检查，代码和格式串则被编译器整个删掉，热路径上没有任何代价。打开后，每一层在事件发生、锁已经放开之后各打印一行，前缀是当前进程的 pid 和名字：

\begin{itemize}
  \item \textbf{page}：\code{alloc\_pages}/\code{kalloc} 的阶、物理地址、来自 pcp 还是伙伴系统；\code{\_\_free\_pages}/\code{kfree}/\code{put\_page} 的释放；pcp 的批量补充、超过 high 后的回吐和 \code{drain\_all\_pages}。
  \item \textbf{slab}：\code{kmem\_cache\_alloc}/\code{free} 的对象地址、所在 slab 及其占用（如 \code{16/25}）、新建 slab、空 slab 还给伙伴系统；\code{kmalloc} 选中的尺寸类或大块的阶。
  \item \textbf{walk}：\code{walk()} 为某个虚拟地址新建的 L2/L3 页表页；每次用户缺页处理完后，该地址经过的 L1/L2/L3 三个描述符和最终物理地址。
  \item \textbf{fault} 与 \textbf{cow}：缺页地址、读写、是否来自内核的 copyin/copyout，以及结果——写时复制拷贝或原地复用、brk 堆的零页、VMA 缺页，或 OOM/SIGSEGV/SIGBUS。
  \item \textbf{brk}、\textbf{mmap}、\textbf{munmap}：系统调用的参数和返回值。
\end{itemize}

页、slab 和 kmalloc 的行以调用者的返回地址结尾（\code{\_\_builtin\_return\_address(0)}）。内核里没有符号表，主机端的 \file{tools/mm\_trace\_sym.py} 用同一次构建的 \file{kernel/kernel.sym} 把它换成“函数名$+$偏移”。下面是 shell 执行 \code{echo hi} 时的一段（地址略去高位）：

\begin{txtcode}
mm[16 sh] slab: proc_cache alloc ...8029780 (slab pa 0x8029000 6/10) <- allocproc+0x40
mm[16 sh] page: alloc order 0 pa 0x3b638000 (pcp) <- alloc_kstack+0x18
mm[16 sh] page: alloc order 0 pa 0x3b636000 (pcp) <- walk+0xd8
mm[16 sh] walk: va 0x0: new L2 table pa 0x3b636000 in L1[0]
mm[16 sh] slab: vma_cache alloc ...80310c0 (slab pa 0x8031000 3/42) <- dup_mmap+0x5c
mm[16 sh] page: alloc order 0 pa 0x3b632000 (pcp) <- cow_break+0x134
mm[16 sh] cow: va 0xa000 copied pa 0x3b641000 -> new pa 0x3b632000
mm[17 sh] page: put_page -> free order 0 pa 0x3b641000 <- cow_break+0x1ac
mm[17 sh] cow: va 0xa000 copied pa 0x3b641000 -> new pa 0x3b61f000
mm[17 sh] brk: sbrk(131072) 0xb000 -> 0x2b000 (pages are faulted in lazily)
mm[17 sh] page: alloc order 0 pa 0x3b61e000 (pcp) <- handle_user_mm_fault+0x180
mm[17 sh] fault: va 0xb008 write -> brk heap, mapped
mm[17 sh] walk: va 0xb008  L1[0]=0x3b636003 L2[0]=0x3b635003 L3[11]=0x6000003b61e74b
                -> pa 0x3b61e008
mm[17 sh] page: alloc order 0 pa 0x3b619000 (pcp) <- e2_file_read_at+0xb0
mm[17 sh] page: free  order 0 pa 0x3b619000 <- e2_file_read_at+0x128
\end{txtcode}

一次 fork 的全过程在这里一目了然：\code{allocproc} 从 slab 拿 \code{struct proc}，从 pcp 拿内核栈，\code{walk} 为子进程建页表，\code{dup\_mmap} 复制 VMA；随后父子进程几乎同时写同一页（\code{0xa000}），两边检查时引用计数都还是 2，于是各自复制，后放手的那次 \code{put\_page} 把旧页的引用降到 0、交还分配器；子进程的 \code{sbrk} 只挪动 brk，真正的页在第一次写时才由缺页处理分配。L3 描述符 \code{0x6000003b61e74b} 也可以对照第~\ref{sec:walk-mappages} 节逐位读出：UXN/PXN 置位（数据不可执行）、AF、Inner Shareable、EL0 可读写、Normal 内存。

跟踪对时序影响很大——串口每行要花几毫秒——所以它适合在单个命令上打开，用 \code{MM\_TRACE\_PID} 和 \code{MM\_TRACE\_LIMIT} 收窄范围；跑 usertests 这类压力测试时应保持关闭。

\section{最终内核页表}

\subsection{\code{kvmmake()}}

\begin{ccode}
pagetable_t kvmmake(void) {
  pagetable_t kpgtbl = (pagetable_t) kalloc();
  memset(kpgtbl, 0, PGSIZE);
  // BCM2837 外设窗口 16 MiB，以及 ARM-local 外设一页：设备内存，不可执行
  kvmmap(kpgtbl, PERIPHERAL_BASE, PERIPHERAL_BASE_PA, 0x01000000, PTE_DEVICE | PTE_XN);
  kvmmap(kpgtbl, LOCAL_BASE, LOCAL_BASE_PA, PGSIZE, PTE_DEVICE | PTE_XN);
  kvmmap(kpgtbl, RAMDISK, RAMDISK_PA, RAMDISK_SIZE, PTE_NORMAL | PTE_XN);
  // 摄像头 DMA：不可缓存的普通内存
  kvmmap(kpgtbl, CAMDMA, CAMDMA_PA, CAMDMA_SIZE, PTE_NORMAL_NC | PTE_XN);
  // 可装载模块区：.ko 的代码拷进来后要执行，所以这里不加 PTE_XN
  kvmmap(kpgtbl, MODAREA, MODAREA_PA, MODAREA_SIZE, PTE_NORMAL);
  // 物理第 0 页：spin-table 槽 0xe0..0xf0，唤醒副核时要写
  kvmmap(kpgtbl, KERNBASE, 0, PGSIZE, PTE_NORMAL | PTE_XN);
  // 内核代码：只读、可执行
  kvmmap(kpgtbl, KERNLINK, V2P(KERNLINK), (uint64)etext-KERNLINK, PTE_NORMAL | PTE_RO);
  // 内核数据和启动窗口其余部分：4 KiB 页，可读写、不可执行
  kvmmap(kpgtbl, (uint64)etext, V2P(etext), (uint64)P2V(EARLYTOP)-(uint64)etext,
         PTE_NORMAL | PTE_XN);
  // memblock 知道的其余 RAM：尽量用 2 MiB 块（Linux 的 map_mem()）
  for(int i = 0; i < memblock.memory.cnt; i++){
    uint64 start = max(memblock.memory.regions[i].base, EARLYTOP);
    uint64 stop  = memblock.memory.regions[i].base + memblock.memory.regions[i].size;
    // 跳过 [CARVEOUT_START, CARVEOUT_END)，它已按自己的属性映射
    kvmmap_linear(kpgtbl, start, stop, PTE_NORMAL | PTE_XN);
  }
  return kpgtbl;
}
\end{ccode}

代码段只读可执行，数据段可写不可执行（W\^{}X），这样即使内核出现缓冲区溢出，也很难把数据当指令执行。唯一的例外是模块区 \code{MODAREA}（物理 \code{0x07400000} 起）：模块装载器把 \code{.ko} 的代码拷进来再执行，它只能可写又可执行（第~\ref{chap:hotplug} 章）。

\subsection{\code{walk()} 与 \code{mappages()}}
\label{sec:walk-mappages}

所有映射最终都落到这两个函数上。\code{walk()} 从 L1 走到 L3，遇到缺失的中间页表就（按需）分配一页：

\begin{ccode}
pte_t *walk(pagetable_t pagetable, uint64 va, int alloc) {
  if(va >= MAXVA) panic("walk");
  for(int level = 1; level < 3; level++) {
    pte_t *pte = &pagetable[PX(level, va)];
    if((*pte & PTE_VALID) && (*pte & PTE_TABLE)) {
      pagetable = (pagetable_t)P2V(PTE2PA(*pte));
    } else {
      if(!alloc || (pagetable = (pde_t*)kalloc()) == 0)
        return 0;
      memset(pagetable, 0, PGSIZE);
      *pte = PA2PTE(V2P(pagetable)) | PTE_TABLE | PTE_VALID;
    }
  }
  return &pagetable[PX(3, va)];
}

int mappages(pagetable_t pagetable, uint64 va, uint64 size, uint64 pa, uint64 perm) {
  uint64 a = PGROUNDDOWN(va), last = PGROUNDDOWN(va + size - 1);
  for(;;){
    pte_t *pte = walk(pagetable, a, 1);
    if(pte == 0) return -1;
    if(*pte & PTE_AF) panic("mappages: remap");
    *pte = PA2PTE(pa) | perm | PTE_AF | PTE_V;
    if(a == last) break;
    a += PGSIZE; pa += PGSIZE;
  }
  return 0;
}
\end{ccode}

注意页表项里存的是\textbf{物理地址}（硬件遍历页表时用物理地址），而内核软件读写页表时要先 \code{P2V()} 转成虚拟地址。页表页本身也位于被直接映射的 RAM 中，所以内核总能通过 \code{KERNBASE + PA} 访问任何一张页表——这不是递归页表技巧，只是直接映射的自然结果。

\subsection{页表的规模}

用户命令 \code{vmmap} 会调用 \code{kvmdump()} 打印内核页表，每个 2\,MiB 窗口汇总一行，再逐个打印各进程的用户页表。表~\ref{tab:kpgt} 是一次快照的统计（当时还保留着 ramdisk 映射）。

\begin{table}[htbp]
  \centering
  \caption{最终内核页表的组成（一次运行时快照）}
  \label{tab:kpgt}
  \small
  \begin{tabularx}{\linewidth}{@{}lllX@{}}
    \toprule
    L1 项 & 虚拟地址 & L2 有效项 & 内容 \\
    \midrule
    \code{L1[0]} & \code{KERNBASE}$+$0～1\,GiB & 65 & 内核和 RAM（56 项）、ramdisk（1 项）、BCM2837 MMIO（8 项） \\
    \code{L1[1]} & \code{KERNBASE}$+$1～2\,GiB & 1 & ARM-local 外设 \code{0x40000000} \\
    \code{L1[255]} & 接近 \code{MAXVA} & 1 & 进程内核栈及其间的保护页 \\
    \midrule
    \multicolumn{4}{@{}l}{合计：L1 一张、L2 三张、L3 67 张，共 71 页 $=$ 284\,KiB；有效 L3 叶子 32\,961 个} \\
    \bottomrule
  \end{tabularx}
\end{table}

表中的快照来自全部用 4\,KiB 页的旧版本：112\,MiB 就要 56 张 L3 表；照此映射 960\,MiB 需要约 480 张、近 2\,MiB，而且必须在 \code{EARLYTOP} 以下分配，启动页表根本放不下。所以当前版本把 4\,MiB 以上的 RAM 交给 \code{kvmmap\_linear()}：起止地址都 2\,MiB 对齐的部分由 \code{kvmmap\_block()} 直接在 L2 写一个块描述符（\code{PTE\_VALID} 置位、\code{PTE\_TABLE} 清零，输出地址的低 21 位为 0），只有首尾不对齐的零头才落到 4\,KiB 页。960\,MiB 的 RAM 只占 L2 的约 475 项，启动时整张内核页表一共 17 页（见上面 memblock 的 reserved \code{[1]}）。

细粒度属性仍然保得住：代码段只读、数据段不可执行都发生在前 4\,MiB，那里依旧是 4\,KiB 页；carve-out 窗口里的 \code{CAMDMA} 是不可缓存的、\code{MODAREA} 是可执行的，线性映射就绕开整个窗口。代价是 \code{walk()} 现在可能在 L2 遇到块描述符：查询（\code{alloc}$=0$）时返回 0，要求分配时则 \code{panic("walk: block")}——内核不会在线性映射里再去改单页的属性。

\section{用户地址空间}
\label{sec:uvm}

\subsection{内核空间与用户空间}

每个进程有自己的 L1 页表，进程切换时写入 \reg{TTBR0\_EL1}：

\begin{ccode}
void switchuvm(struct proc *p) {
  w_ttbr0_el1(V2P(p->vm->pagetable));
  flush_tlb();
}
\end{ccode}

\subsection{\code{exec()} 之后的布局}

\code{exec()} 建立的用户地址空间从虚拟地址 0 开始（图~\ref{fig:uas}）：先是 ELF 的代码和数据，然后是一页不可访问的保护页和 4 页（16\,KiB）用户栈，栈之上是 \code{sbrk()} 扩展的堆。用户页面的属性是 \code{PTE\_NORMAL | PTE\_U}，代码页可执行。动态链接的程序在主镜像之后还有一页多的 \file{/lib/ld.so}；\code{mmap()} 的区域则远在其上：\code{[MMAP\_MIN, MMAP\_TOP)}，从 64\,GiB（\code{0x10\_0000\_0000}）到 256\,GiB$-$256\,MiB，自顶向下分配，只登记 VMA、按需调页（第~\ref{sec:mmap} 节）。

\begin{figure}[htbp]
  \centering
  \begin{tikzpicture}[x=1cm,y=6mm]
    \foreach \y/\h/\t/\c in {0/1.4/{text / data / bss}/structurecolor!15, 1.4/0.8/{保护页（\code{uvmclear} 清除 U 位）}/gray!25,
        2.2/1.6/{用户栈 16\,KiB（\code{USTACKPAGES}=4）}/structurecolor!25, 3.8/1.2/{堆（\code{sbrk}）}/structurecolor!8}{
      \draw[fill=\c] (0,\y) rectangle (6,\y+\h);
      \node[font=\small] at (3,\y+\h/2) {\t};
    }
    \node[lab, left] at (0,0) {VA 0};
    \node[lab, left] at (0,5) {\code{p->vm->sz}};
    \draw[arr] (6.3,3.8) -- (6.3,2.4) node[midway, right, lab]{栈向下增长};
  \end{tikzpicture}
  \caption{\code{exec()} 之后的用户地址空间}
  \label{fig:uas}
\end{figure}

进程的内核栈不在用户页表里，而在内核页表高端的 \code{KSTACK(slot)} 处，每个栈一页，相邻栈之间隔一页不映射的保护页。内核栈溢出会立刻触发翻译错误，而不是悄悄覆盖相邻进程的栈。第~\ref{chap:process} 章会介绍内核栈的动态分配。

\subsection{第一个用户进程的页表}

所有用户地址空间都由同一组函数建立。\code{uvmcreate()} 分配一页并清零，作为空的 L1 页表——全 0 意味着全部表项无效，这一点很重要：任何从 \code{kalloc()} 拿来做页表的页都必须先清零，否则 \code{0x05} 填充的垃圾会被 MMU 当成有效描述符。第一个进程没有 ELF 文件可读，内核把一小段编进内核镜像的 \code{initcode} 直接拷进一页：

\begin{ccode}
void uvminit(pagetable_t pagetable, uchar *src, uint sz) {
  if(sz >= PGSIZE) panic("inituvm: more than a page");
  char *mem = kalloc();
  memset(mem, 0, PGSIZE);
  mappages(pagetable, 0, PGSIZE, V2P(mem), PTE_NORMAL|PTE_U);
  memmove(mem, src, sz);              // 通过内核直接映射写入
  uvmsync_icache(pagetable, sz);      // 写入的是指令，要同步 I-cache
}
\end{ccode}

\code{userinit()} 再把 \code{trapframe} 的 \code{elr} 设为 0、\code{sp} 设为 \code{PGSIZE}：进程第一次“返回”用户态时，从虚拟地址 0 取第一条指令，栈从这一页的顶部向下长。同一个物理页此时有两个虚拟地址：内核通过 \code{KERNBASE+PA} 写它，用户通过 VA 0 执行它。拷贝时用的是前者，所以内核根本不需要先切换到这个进程的页表。

\subsection{\code{fork}：写时复制}
\label{sec:cow-fork}

\code{fork()} 要给子进程一份与父进程内容相同、但互不影响的内存。最直接的做法是对 \code{[0, sz)} 的每一页 \code{kalloc()} 一页、拷贝、再映射进子进程——早期的 \code{uvmcopy()} 就是这样。代价是 fork 一个 8\,MiB 的进程就要复制 8\,MiB，而 shell 的子进程往往紧接着 \code{exec()}，复制的内容立刻被丢弃。vfork 能解决 fork+exec 的场合（第~\ref{sec:clone-flags} 节），但它要求子进程在 \code{exec} 之前什么都不改。

有了第~\ref{sec:buddy} 节的 \code{struct page} 和它的 \code{refcount}，当前源码树改用写时复制（copy-on-write，COW）：fork 时\textbf{只共享、不复制}，谁先写谁复制。

\paragraph{fork 时：共享并写保护.} \code{uvmcopy()} 把父进程的每个私有页以同一物理地址映射进子进程，并给这一页加一份引用；原来用户可写的页，在\emph{父子两边}都改成“只读 + \code{PTE\_COW}”：

\begin{ccode}
if(*pte & PTE_SHARED){                    // libc.so 代码：归缓存所有，照旧共享
  mappages(new, i, PGSIZE, pa, flags | PTE_SHARED);
  continue;
}
if((*pte & PTE_U) && (*pte & PTE_RO) == 0)
  *pte |= PTE_RO | PTE_COW;               // 父进程：可写 -> COW
if(*pte & PTE_COW)
  flags |= PTE_RO | PTE_COW;
mappages(new, i, PGSIZE, pa, flags);      // 子进程：同一物理页
get_page(virt_to_page(P2V(pa)));          // refcount++
\end{ccode}

\code{PTE\_COW} 是描述符里留给软件的第 57 位（MMU 忽略 \code{[58:55]}，第 55、56 位已被 \code{PTE\_SHARED}、\code{PTE\_SPECIAL} 占用），意思是“只读只是因为共享，写了就复制”。本来就只读的私有页（代码段、只读数据、栈的保护页）不加 \code{PTE\_COW}，只增加引用、两边共享，写它们仍是致命错误。\code{dup\_mmap()} 对 \code{MAP\_PRIVATE} 的 VMA 页做同样的处理，共享映射和设备帧则照旧。复制不再写入新字节，所以原来跟在 \code{uvmcopy()} 后面的 \code{uvmsync\_icache()} 也删掉了。

整个过程在父进程的 \code{vm->lock} 下进行；结束后一次 \code{flush\_tlb()}（\code{TLBI VMALLE1IS}）广播作废，父进程其他线程在别的 CPU 上缓存的“可写”翻译随之失效。

\paragraph{写的时候：\code{cow\_break()}.} 写一个 COW 页会产生权限错误（\reg{ESR\_EL1} 的 EC $=$ \code{0x24}、WnR $=1$、DFSC $=$ \code{0b0011xx}）。\code{usertrap()} 照常把它交给 \code{handle\_user\_mm\_fault()}，后者在按地址区分 brk 堆和 VMA 之前，先看这是不是一个 COW 页——代码段、数据段和栈不属于任何 VMA，不先判断就会被当成越界：

\begin{ccode}
pa = PTE2PA(*pte);
pg = virt_to_page(P2V(pa));
if(page_count(pg) == 1){                  // 别人都已复制或退出：就地收归己有
  *pte &= ~(PTE_RO | PTE_COW);
  flush_tlb();
  return 1;
}
mem = kalloc();                           // 复制一份
memmove(mem, P2V(pa), PGSIZE);
if((*pte & PTE_UXN) == 0)                 // 可执行页还要同步 I-cache
  sync_icache(mem);
*pte = PA2PTE(V2P(mem)) | (*pte & ~(PA2PTE(~0ULL) | PTE_RO | PTE_COW));
flush_tlb();                              // 先让所有 CPU 忘掉旧翻译
put_page(pg);                             // 再放弃对共享页的那一份引用
\end{ccode}

顺序很重要：必须先广播作废 TLB 再 \code{put\_page()}。另一方可能恰好是最后一个使用者，引用一减到 0，伙伴系统就会把这页发给别人；如果此时某个 CPU 上还留着旧翻译，就会写进别人的内存。

返回用户态后，那条 \code{str} 重新执行，这次写进自己的页。\code{cow\_break()} 不会睡眠（\code{kalloc()} 只用自旋锁），这一点后面会用到。

\paragraph{释放：\code{kfree} 换成 \code{put\_page}.} 有了共享，“释放一页”的含义变成“放弃我这一份”。\code{uvmunmap()}（进程退出、\code{exec} 换镜像）、\code{growproc()} 缩小堆、\code{zap\_range()} 撤销私有 VMA，这几处原来的 \code{kfree()} 都改成 \code{put\_page(virt\_to\_page(...))}，最后一个使用者离开时页才回到伙伴系统。第~\ref{sec:buddy} 节的检查在这里派上了用场：如果哪条路径漏改，仍对共享页调用 \code{kfree()}，会立刻因 “still has 2 references” 而 \code{panic}，而不是悄悄地把别的进程还在用的页放回空闲链表。

\paragraph{\code{copyout()} 也要遵守.} 内核写用户内存走的是自己的直接映射，用户 PTE 的 AP 位对它不起作用。如果不处理，子进程 \code{read(fd, buf, n)} 时内核会把数据直接写进仍与父进程共享的物理页，父进程的内存就被改了。\code{copyout()} 原来就会对只读页调用 \code{uva2ka\_writable()} 拒绝写入，现在在拒绝之前先试 \code{cow\_break()}。而且对当前地址空间，检查、破除 COW 和这一页的 \code{memmove} 都在 \code{vm->lock} 下完成：fork 也持这把锁，于是另一个线程的 \code{fork()} 不可能恰好在检查和写入之间把这一页重新变成共享页。因为 \code{cow\_break()} 不睡眠，持有自旋锁的调用者（比如在锁内回写退出状态的 \code{wait()}）也能正确写入 COW 页；只有需要按需调页的情形才在锁外走 \code{uva2ka\_fault()}，然后在锁内重试这一页。

\paragraph{统计与测试.} \file{/proc/meminfo} 多了三行：\code{CowShared}（fork 时共享出去的页数）、\code{CowCopied}（写时复制的次数）、\code{CowReused}（最后一个使用者就地收回的次数）。新的用户程序 \code{cowtest} 覆盖五种情形：

\begin{itemize}
  \item \textbf{big}：先 \code{sbrk} 占用三分之二的空闲内存，再 fork——不共享的话需要两倍内存，必然失败；
  \item \textbf{isolate}：父子写同一批页，各自只看到自己的修改；
  \item \textbf{three}：两个子进程先后复制并退出，父进程最后写时 \code{refcount} 已是 1，走就地收回的路径；
  \item \textbf{copyout}：子进程 \code{read()} 一个管道到 COW 页，父进程的页不变；
  \item \textbf{mmap}：\code{MAP\_PRIVATE|MAP\_ANONYMOUS} 页同样写时复制。
\end{itemize}

全部测试在一个子进程里运行，子进程退出后检查空闲页数回到原值：

\begin{shcode}
$ cowtest
PASS big: sbrk two thirds of free memory
PASS big: fork succeeds (would need 2x memory without COW)
PASS big: child reads parent's data
PASS big: pages were shared
...
PASS mmap: parent's MAP_PRIVATE pages unchanged
MemFree before 110524 kB, after 110524 kB
cowtest: OK (0 failures)
\end{shcode}

\paragraph{行为上的变化.} COW 把“内存不够”从 fork 时推迟到了写的时候：fork 几乎总能成功，真正复制时如果没有空闲页，\code{cow\_break()} 返回 OOM，进程被杀（\code{out of memory at va=...}）——这与 Linux 的 OOM 行为一致。另一个可见的变化在 usertests 的 \code{forktest}：本内核没有固定大小的进程表，原来 fork 不到 1000 次就会因内存耗尽而失败，测试据此判定“fork 能正确报告失败”。有了 COW，每个子进程只花几页（内核栈、trapframe、页表页），fork 一直能成功到三千多个子进程，测试反而报错；更糟的是，最终耗尽内存的往往是父进程自己写 COW 页时的那次复制，于是父进程被 OOM 杀掉，几千个 \code{struct proc}、\code{vmspace} 所在的 slab 页也被 \code{kmem\_cache} 永久留作缓存（表~\ref{tab:buddy-vs-linux} 中 slab 一行的“无收缩”）。所以 \file{param.h} 加了任务数上限 \code{NPROC\_MAX}，取 Linux \code{threads-max} 的默认算法“内存页数除以 32”，112\,MiB 对应 896；\code{newprocslot()} 在已经建立 896 个进程描述符时直接返回失败，\code{fork()} 返回 $-1$，相当于 Linux 的 \code{EAGAIN}。描述符在进程回收后会被复用，所以这个上限约束的是同时存在的任务数。

\subsection{TLB 与地址空间切换}

如果每次访存都要走三级页表，一次 \code{ldr} 就会变成四次内存访问。MMU 因此把最近用过的翻译结果缓存在 TLB（Translation Lookaside Buffer）里，命中时不再遍历页表。代价是：一旦软件改了页表，TLB 里可能还留着旧的翻译，必须由软件显式作废。

进程切换正是最大的一次“改页表”——\reg{TTBR0} 换成另一棵树，之前缓存的低地址翻译全都过期了：

\begin{ccode}
void switchuvm(struct proc *p) {
  if(p->vm == 0 || p->vm->pagetable == 0) panic("switchuvm: no pagetable");
  w_ttbr0_el1(V2P(p->vm->pagetable));
  flush_tlb();                 // dsb ishst; tlbi vmalle1is; dsb ish; isb
}
\end{ccode}

\code{flush\_tlb()} 作废全部 EL1 翻译，连内核的也一起丢掉了。ARMv8 提供了 ASID（地址空间标识）：把一个 8 或 16 位的进程编号写进 \reg{TTBR0} 的高位，TLB 表项会带上这个标签，切换时就不必全部作废。xv6-aarch64 没有使用 ASID，每次切换都整体作废，换来的是实现简单、不必管理 ASID 的回收；对于这个规模的系统，多出来的 TLB 缺失是可以接受的。注意内核页表（\reg{TTBR1}）在整个运行期间只装一次（\code{kvminithart()}），这也是把两个地址空间分开的好处。

\subsection{缺页与按需分配}

如果用户进程访问了一个没有映射的地址，页表遍历在某一级遇到无效项，CPU 产生同步异常进入 \code{usertrap()}。此时 \reg{ESR\_EL1} 的 EC 字段是 \code{0x24}（来自 EL0 的数据中止）或 \code{0x20}（来自 EL0 的指令中止）；ISS 的低 6 位 DFSC 说明原因，\code{0b0001xx} 是第 \code{xx} 级的翻译错误，\code{0b0011xx} 是权限错误；bit 6（WnR）表示是写；出错的虚拟地址在 \reg{FAR\_EL1} 里。

很多内核利用这个异常实现\emph{按需分配}（demand paging）：进程第一次碰到某页时触发翻译错误，内核检查 \reg{FAR\_EL1} 落在合法区间里，就当场分配、映射，然后 \code{eret} 重新执行那条出错的指令。进程看到的只是一条稍微慢一点的 \code{ldr}。

xv6-aarch64 对两段地址采用不同的策略：

\begin{itemize}
  \item \code{[0, sz)}：\code{exec()} 和 \code{growproc()} 调用 \code{uvmalloc()} 立即分配并清零所有页，合法访问永远不会缺页；
  \item mmap 区：\code{mmap()} 只登记 VMA，第一次访问时由 \code{user\_page\_fault()} $\to$ \code{handle\_mm\_fault()} 按 VMA 分配或读入（第~\ref{sec:mmap} 节）。共享文件页的第一次写也走这里，借权限错误标记脏页。
\end{itemize}

其余异常——访问 \code{[0, sz)} 之外又不在任何 VMA 里的地址、写只读页、执行不可执行页——一律打印 \code{esr}、\code{elr}、\code{far}，调用 \code{uvmdump()} 打印这个进程的整张页表（mmap 区里的错误只打印一行 SIGSEGV/SIGBUS），然后杀掉进程：

\begin{ccode}
} else if((ec == 0x24 || ec == 0x20) && user_page_fault(p, esr)){
  // mmap 区的缺页已处理，或已打印 SIGSEGV/SIGBUS 并设 killed
} else {
  printf("usertrap(): unexpected esr=%p ec=%p pid=%d\n", esr, ec, p->pid);
  printf("            elr=%p far=%p\n", r_elr_el1(), r_far_el1());
  uvmdump(p->vm->pagetable, p->pid, p->name, "usertrap");
  p->killed = 1;
}
\end{ccode}

这份输出正是 \ref{sec:caches} 节分析真机故障时用的线索。

\section{\code{mmap()}：VMA、按需调页与设备映射}
\label{sec:mmap}
\label{sec:pcache-mmap}

前面的用户地址空间 \code{[0, sz)} 由 \code{exec()} 和 \code{sbrk()} 一次性建好，页表里有什么，进程就有什么。\code{mmap()} 换了一种思路，仿照 Linux：先用一个 \code{vm\_area\_struct}（VMA）\emph{描述}一段地址\emph{应该}装着什么——哪个文件的哪一段、是否与别的进程共享、是否是设备的物理帧——页表只记录已经碰过的页。第一次访问某页时 MMU 报缺页，内核按 VMA 的描述把这一页“变出来”，再让那条指令重新执行。图~\ref{fig:mmap-structs} 是数据结构，图~\ref{fig:mmap-paths} 是三条主要路径。

\begin{figure}[htbp]
  \centering
  \includegraphics[width=\linewidth]{fig-y-mmap-structs.pdf}
  \caption{mmap 的数据结构：地址空间布局、\code{vmspace}/\code{vm\_area\_struct}、四类物理页的归属与 PTE 位}
  \label{fig:mmap-structs}
\end{figure}

\subsection{\code{vm\_area\_struct} 与地址布局}

\code{struct vmspace}（相当于 Linux 的 \code{mm\_struct}）多了两个成员：按地址排序的 VMA 链表 \code{mmap}，以及一把可睡眠的读写信号量 \code{mmap\_lock}。缺页处理拿读锁，\code{mmap}/\code{munmap} 拿写锁；改页表本身仍由自旋锁 \code{vm->lock} 保护。VMA 定义在 \file{include/kernel/mm.h}，字段是 Linux 的一个子集，\code{vm\_flags} 的位值也照抄 Linux：

\begin{ccode}
struct vm_area_struct {
  uint64 vm_start, vm_end;    // [start, end)，页对齐
  uint64 vm_flags;            // VM_READ/WRITE/EXEC/SHARED/MAYWRITE/PFNMAP/IO...
  uint64 vm_page_prot;        // 内存类型：PTE_NORMAL / PTE_NORMAL_NC / PTE_DEVICE
  uint64 vm_pgoff;            // 文件偏移（或物理帧号），以页为单位
  struct file *vm_file;       // 后备文件，匿名映射为 0
  const struct vm_operations_struct *vm_ops;   // open/close/fault
  void *vm_private_data;
  uint64 vm_shmkey;           // 共享匿名内存的页缓存键
  struct vmspace *vm_mm;
  struct vm_area_struct *vm_next;
};
\end{ccode}

mmap 区在 \code{[MMAP\_MIN, MMAP\_TOP)}，即 64\,GiB 到 256\,GiB$-$256\,MiB，不带 \code{MAP\_FIXED} 时 \code{get\_unmapped\_area()} 自顶向下找空洞；\code{growproc()} 不允许堆长进最低的 VMA。\code{do\_mmap()} 检查参数、分配 VMA，然后分两种情况：普通文件和匿名映射到此为止，一页都不分配；字符设备则调用驱动的 \code{f\_op->mmap(f, vma)}，签名与 Linux 相同。

\subsection{按需调页：\code{handle\_mm\_fault()}}

\begin{figure}[htbp]
  \centering
  \includegraphics[width=\linewidth]{fig-z-mmap-paths.pdf}
  \caption{mmap 的三条路径：(a) 缺页，(b) \code{munmap}/\code{msync} 回写，(c) 设备 mmap}
  \label{fig:mmap-paths}
\end{figure}

\code{usertrap()} 遇到 EC $=$ \code{0x24}（数据中止）或 \code{0x20}（指令中止）时先交给 \code{user\_page\_fault()}。地址低于最低 VMA 的缺页仍走原来的详细报告；否则打开中断（读文件可能睡眠）调用 \code{handle\_mm\_fault()}：

\begin{ccode}
if(fsc < 0x04 || fsc > 0x0f || far < mmap_lowest(p->vm))
  return 0;                                   // 不是 mmap 区：保留原来的崩溃报告
if(ec == 0x20)                     flags |= FAULT_FLAG_INSTR;
else if((iss & (1<<6)) && !(iss & (1<<8))) flags |= FAULT_FLAG_WRITE;  // WnR
intr_on();
r = handle_mm_fault(p->vm, far, flags);       // 失败：打印 SIGSEGV/SIGBUS 并杀进程
\end{ccode}

\code{handle\_mm\_fault()} 在 \code{mmap\_lock} 读锁下找到 VMA，检查权限（不在任何 VMA 里、写只读区、执行不可执行区都是 SIGSEGV），然后按四种情况产生物理页：

\begin{enumerate}
  \item \textbf{PTE 已存在}：只可能是“写一个干净的共享文件页”，把 AP 改成可写即可（见下文的脏页跟踪）；
  \item \textbf{\code{VM\_PFNMAP}}（设备）：由 \code{vm\_ops->fault} 给出物理帧号；
  \item \textbf{\code{VM\_SHARED}}：到共享页缓存里找或建这一页；
  \item \textbf{私有}：\code{kalloc()} 一页，文件映射用 \code{f\_op->read\_at} 读入（整页都在文件末尾之后则 SIGBUS，尾部补 0），匿名映射清零。
\end{enumerate}

最后在 \code{vm->lock} 下 \code{walk()} 并写入 PTE。如果另一个线程抢先装好了同一页，就把自己的页退回；然后返回，\code{eret} 重新执行那条 \code{ldr}/\code{str}。

\subsection{共享页：\code{pcache\_page} 与 \code{shm\_object}}

\code{MAP\_SHARED} 要求所有映射同一文件同一页的 PTE 指向同一个物理页，内核因此需要一个以“对象 $+$ 页号”为键的表：

\begin{ccode}
struct pcache_page {
  struct pcache_page *next;   // 64 桶哈希链
  uint64 key;                 // inode* 或 shm_object*
  uint64 index;               // 页号
  uint64 pa;
  int mapcount;               // 有多少个 PTE 指向它
};
\end{ccode}

文件页的键是 inode 指针，所以两个进程各自 \code{open()} 同一文件再 \code{mmap}，也会共享同一页；\code{mapcount} 降到 0 时页被释放。共享匿名内存（\code{MAP\_SHARED|MAP\_ANONYMOUS} 和 \code{/dev/zero}）没有后备文件，键是一个按 VMA 计数的 \code{shm\_object}，页要活到最后一个 VMA 消失——否则 fork 出的子进程先写、随后退出，那一页就跟着丢了（\code{mmaptest} 抓到过这个错误）。

\paragraph{脏页跟踪.} Cortex-A53 是 ARMv8.0，没有硬件脏位（DBM 是 ARMv8.1 的特性）。所以共享文件页第一次总是装成\emph{只读}；第一次写触发权限错误，\code{handle\_mm\_fault()} 的情况 1 把 AP 改成可写并刷 TLB。于是“共享文件 VMA 里可写的 PTE”就等于“脏页”。

\paragraph{回写.} \code{munmap()}、进程退出和 \code{exec()} 都经过 \code{zap\_range()}：先对可写 PTE 指向的页调用 \code{file\_write\_at()} 写回文件（只写到 \code{i\_size}，不扩展文件），这一步会睡眠，所以不持 \code{vm->lock}；然后在 \code{vm->lock} 下每批 32 个 PTE 清零、\code{flush\_tlb()}，最后才释放物理页。\code{msync(MS\_SYNC)} 顺序相反：先把脏 PTE 改回只读并刷 TLB，再写回——写回期间的并发写会重新缺页、重新变脏，不会丢。

\subsection{设备 mmap：\code{f\_op->mmap} 与 \code{remap\_pfn\_range()}}

设备没有页缓存，驱动知道的是物理地址。驱动的 \code{mmap} 回调检查偏移和长度后调用 \code{remap\_pfn\_range()} 立即装好整段 PTE，例如帧缓冲：

\begin{ccode}
static int fb_mmap(struct file *f, struct vm_area_struct *vma) {
  uint64 size = vma->vm_end - vma->vm_start, len = PGROUNDUP(fb.size);
  if(!fb.ready || vma->vm_pgoff > len / PGSIZE ||
     size > len - vma->vm_pgoff * PGSIZE || (vma->vm_flags & VM_EXEC))
    return -1;
  return remap_pfn_range(vma, vma->vm_start,
                         (fb.base >> PGSHIFT) + vma->vm_pgoff, size, PTE_NORMAL_NC);
}
\end{ccode}

\code{remap\_pfn\_range()} 给 VMA 加上 \code{VM\_IO|VM\_PFNMAP}，并在每个 PTE 上打软件位 \code{PTE\_SPECIAL}（bit 56）：这一页不归 \code{kalloc} 管，\code{munmap} 只清 PTE、不释放，fork 照抄 PTE。内存类型由驱动决定：MMIO 必须是 Device-nGnRnE，帧缓冲和 DMA 缓冲用 Normal-NC（写合并，并与内核对同一物理页的别名属性一致）。

\begin{table}[htbp]
  \centering
  \caption{支持 \code{mmap} 的设备}
  \label{tab:dev-mmap}
  \small
  \begin{tabularx}{\linewidth}{@{}llX@{}}
    \toprule
    设备 & 主/次 & 行为 \\
    \midrule
    \code{/dev/mem} & 8/1 & 偏移即物理地址；拒绝落在 memblock 内存里的区间（carve-out 窗口除外，相当于 \code{STRICT\_DEVMEM}）；MMIO 用 Device-nGnRnE，其余 Normal-NC \\
    \code{/dev/zero} & 8/5 & 变成匿名映射；\code{MAP\_SHARED} 即共享匿名内存 \\
    \code{/dev/fb0} & 9/0 & \file{bcm2708\_fb.c}：首次 open 用一条多 tag 的 mailbox 消息分配 640$\times$480$\times$32 帧缓冲；支持 \code{FBIOGET\_VSCREENINFO}/\code{FSCREENINFO} \\
    \code{/dev/video0} & 5 & Unicam DMA 缓冲；地址由内核选；\code{VM\_NOFILEREF} 让映射不持有 file，最后一个 fd 关闭即停流；\code{vm\_ops} 的 open/close 维护 \code{camera\_mapped} 计数 \\
    \bottomrule
  \end{tabularx}
\end{table}

\subsection{内核替用户访问 mmap 区}

\code{copyin}/\code{copyout} 现在可能碰到还没缺页进来的地址。规则有三条：

\begin{itemize}
  \item \code{read}/\code{write}/\code{pread}/\code{pwrite} 在拿 inode 锁\emph{之前}先 \code{fault\_in\_user\_range()}，把整块用户缓冲区缺页进来；
  \item \code{copyout} 仍遇到缺页时调用 \code{uva2ka\_fault()}：若当前持有自旋锁（例如 \code{pipewrite}、\code{consoleread} 持锁拷贝）直接失败，不睡眠；否则以 \code{FAULT\_FLAG\_KERNEL} 调 \code{handle\_mm\_fault()}，只 \code{down\_read\_trylock(mmap\_lock)}，避免与正在回写的 \code{munmap} 形成 ABBA 死锁；
  \item \code{PTE\_SPECIAL} 且不在内核直接映射范围内的页（帧缓冲、MMIO）一律返回 \code{EFAULT}，内核不去碰没有映射的物理地址。
\end{itemize}

\code{fetchaddr()} 也不再要求地址在 \code{[0, sz)}，只要在 256\,GiB 用户空间内，这样 \code{argv} 可以放在 mmap 区里。

\subsection{fork、vfork、exit 与 \code{exec}}

\code{fork()} 在 \code{allocproc()} 之前拿 \code{mmap\_lock} 读锁，复制完 \code{[0, sz)} 后调用 \code{dup\_mmap()}：每个 VMA 复制一份（\code{filedup}、\code{shm\_get}、\code{vm\_ops->open}），私有页按写时复制共享（第~\ref{sec:cow-fork} 节），共享页 \code{mapcount++}、PTE 照抄，设备帧照抄 PTE；中途失败用不睡眠的 \code{mmap\_abort()} 回滚。带 \code{CLONE\_VM} 的线程和 vfork 子进程直接共享 \code{vmspace}，不复制。

共享 \code{vmspace} 带来一个问题：vfork 子进程退出时不能把父进程的映射拆掉。所以 \code{vmspace} 有两个计数（第~\ref{sec:clone-flags} 节）：\code{refcount} 记录还有谁持有页表，\code{users} 记录还有几个活着的任务在用。\code{exit()}、线程退出和 vfork 子进程的 \code{exec()} 都调用 \code{vmspace\_leave()}，只有 \code{users} 减到 0 时才 \code{exit\_mmap()}。\code{exec()} 对私有 \code{vmspace} 在提交新镜像前 \code{exit\_mmap()}；vfork 子进程则换到 \code{vmspace\_wrap()} 新建的私有 \code{vmspace}，旧的 VMA 原样留给父进程。

\subsection{测试与限制}

\file{user/mmaptest.c} 在 QEMU raspi3b 上 49 项全部通过，连续运行 \file{/proc/meminfo} 的 MemFree 不变。它覆盖：匿名零页、\code{munmap} 中间一页拆分 VMA、访问空洞或只读页被杀、\code{MAP\_FIXED}、fork 后的共享与私有语义、子进程先写的共享页、vfork 子进程写父进程的映射并退出、\code{MAP\_PRIVATE} 文件（关闭 fd 后仍可用、尾部补零、写不回文件）、\code{MAP\_SHARED} 文件（\code{msync}、\code{munmap} 回写、文件大小不变、两个进程各自 open 后共享页）、\code{read()} 写进未触碰的映射、超出 EOF 的 SIGBUS、\code{/dev/zero}、用 \code{/dev/mem} 读系统定时器、\code{/dev/mem} 拒绝内核 RAM、\code{/dev/fb0} 画彩条。

已知限制：没有写时复制，fork 复制所有已触碰的私有页；\code{pcache\_page} 只在页被映射期间存在，不是统一的页缓存，\code{write()} 不会更新已映射的共享页；没有 \code{mprotect}、\code{mremap}、\code{madvise} 和页回收。

\section{MMIO、总线地址与 DMA}

\subsection{MMIO}

SoC 把外设寄存器和 RAM 放进同一个物理地址空间，CPU 用普通的 \code{ldr}/\code{str} 访问它们，由片内地址译码器决定一次访问是去 DDR 控制器还是去某个外设（表~\ref{tab:physmap}）。

\begin{table}[htbp]
  \centering
  \caption{Raspberry Pi 3 的物理地址布局（本内核使用的部分）}
  \label{tab:physmap}
  \small
  \begin{tabular}{@{}ll@{}}
    \toprule
    物理地址 & 用途 \\
    \midrule
    \code{0x00000000}～\code{0x3bffffff} & RAM，以 DTB 的 \code{/memory} 为准（QEMU 960\,MiB；无 DTB 时 112\,MiB） \\
    \code{0x00080000} & 内核镜像装载地址 \\
    \code{0x07200000}～\code{0x072fffff} & 摄像头 DMA 区（\code{CAMDMA}） \\
    \code{0x3f003000} & BCM System Timer \\
    \code{0x3f00b000} & legacy 中断控制器；\code{+0x880} 处为 property mailbox \\
    \code{0x3f101000} & 时钟管理器（摄像头 \code{CM\_CAM1}） \\
    \code{0x3f200000} & GPIO \\
    \code{0x3f202000} & SDHOST（外置 SD 卡） \\
    \code{0x3f205000} & BSC0（I2C） \\
    \code{0x3f215000} & AUX / Mini UART \\
    \code{0x3f300000} & Arasan SDHCI（板载 Wi-Fi） \\
    \code{0x3f801000} & Unicam CSI1 \\
    \code{0x3f980000} & DWC2 USB 主机 \\
    \code{0x40000000} & ARM-local 中断与定时器路由 \\
    \bottomrule
  \end{tabular}
\end{table}

外设寄存器必须映射为 Device 内存。若映射成可缓存的普通内存，写操作可能停留在缓存里迟迟不到达设备，读操作可能拿到旧值，两次写还可能被合并或重排——而对设备来说，“写 FIFO”这个动作本身就有副作用。运行时可以用 \code{cat /proc/iomem} 查看内核登记的全部外设地址区间。

\subsection{三种地址视图}

Broadcom 的文档用 VideoCore 总线地址描述外设，例如 Mini UART 在 \code{0x7e215000}。同一组寄存器在不同观察者眼里地址不同：

\begin{txtcode}
VideoCore 总线地址   0x7e215000     (文档、设备树里的写法)
ARM 物理地址         0x3f215000     (CPU 发出的地址)
内核虚拟地址         0xffffff803f215000  (经 TTBR1 翻译)
\end{txtcode}

DMA 方向的地址转换更容易出错。DWC2 USB、Unicam 这些\emph{总线主设备}自己访问内存时，看到的是 VideoCore 总线地址，必须把 ARM 物理地址加上别名前缀：

\begin{ccode}
#define DWC2_DMA_BUS(a) (0xc0000000U | (uint32)V2P(a))
\end{ccode}

\code{0xC0000000} 别名表示“绕过 VideoCore 的 L2 缓存直接访问 SDRAM”。这不是新的内存，也不是 CPU 虚拟地址，而是同一片 RAM 在 GPU 总线上的另一种编址。

\subsection{PCIe：BAR、DMA 与 IOMMU}

PCIe 把两条方向相反的访问路径放在同一条互连上，因而最容易把地址概念混在一起。CPU 访问设备时走
\emph{入站}路径：操作系统给设备的 BAR 分配一段 CPU 物理地址，CPU 对这段 MMIO 窗口执行
\code{ldr}/\code{str}，Root Complex 再把访问路由成 PCIe 请求。设备主动访问主机内存时走
\emph{出站}路径：设备自己的 DMA 引擎发出 PCIe Memory Read/Write 请求，目标是驱动交给设备的
DMA 地址。BAR 和 DMA 因而不是同一套地址，也不是互相替代的功能。

一块 PCIe 设备通常同时涉及四种地址：

\begin{table}[htbp]
  \centering
  \caption{PCIe 设备涉及的地址空间}
  \label{tab:pcie-addresses}
  \small
  \begin{tabularx}{\linewidth}{@{}lX@{}}
    \toprule
    地址 & 含义 \\
    \midrule
    设备内部地址 & 设备片内 SRAM、FIFO 或寄存器在设备内部互连上的地址，通常只由设备逻辑直接使用 \\
    BAR/MMIO 地址 & CPU 访问设备寄存器或可见设备内存窗口时使用的主机物理地址 \\
    DMA/IOVA 地址 & 设备 DMA 请求中携带的地址；驱动把它写进设备描述符或 DMA 寄存器 \\
    主机物理地址 & 数据在主机 DDR 中实际占用的位置 \\
    \bottomrule
  \end{tabularx}
\end{table}

没有 IOMMU 时，DMA 地址通常就是设备可达的主机物理地址；有 IOMMU 时，设备使用 IOVA，
IOMMU 硬件依据内核建立的页表把它翻译成主机物理地址并检查访问权限。IOMMU 不搬运数据，真正
发起读写、产生 PCIe TLP 的仍是设备内部 DMA 引擎。它的价值是隔离设备、支持虚拟机、让连续
IOVA 映射到离散物理页，以及适配设备的地址位宽限制。

DMA 的“数据源”由方向和设备功能决定，并不固定为设备 SRAM。设备向主机写数据时，来源可以是
设备 SRAM，也可以是接收 FIFO、硬件流水线的实时输出，或者经控制器读出的外部介质数据：网卡
从 PHY/MAC 接收包，NVMe 控制器从 NAND 读数据，采集卡从传感器得到数据。它们可能先暂存在
SRAM/FIFO，也可能边产生边由 DMA 写入主机 DDR：

\begin{txtcode}
设备 SRAM / RX FIFO / 硬件流水线 / 外部输入
        -> 设备 DMA 引擎
        -> PCIe Memory Write TLP
        -> IOMMU（可选：IOVA -> 主机 PA）
        -> 主机 DDR
\end{txtcode}

反方向传输时，主机 DDR 才是数据源。以网卡发送为例，驱动在 DDR 中准备数据包与描述符，把
DMA 地址写入网卡；网卡 DMA 引擎发出 Memory Read，从 DDR 取回数据，送进设备 TX FIFO，最后
由 MAC/PHY 发出：

\begin{txtcode}
主机 DDR
        -> IOMMU（可选）
        -> PCIe Memory Read completion
        -> 设备 DMA 引擎
        -> 设备 SRAM / TX FIFO / 处理单元
\end{txtcode}

所以“CPU 通过 BAR 读设备 SRAM”和“设备通过 DMA 把数据写进 DDR”是两件不同的事；即便同一块
设备 SRAM 同时参与二者，前者经过 MMIO 窗口，后者由设备 DMA 引擎完成。PCIe 设备也不一定必须
实现 DMA：低速控制设备可以只暴露 BAR；网卡、NVMe、GPU 等高吞吐设备则必须依靠 DMA，避免 CPU
逐字节搬运。

\section{缓存与 TLB：三个真机故障}
\label{sec:caches}

QEMU 不模拟缓存，也几乎不检查内存属性，下面三个问题都只在真机上出现。它们说明：页表不只是“虚拟地址到物理地址”的映射，还决定了 CPU 怎样访问这块内存。

\subsection{故障一：停在 \code{BPH0abde}}
\label{sec:fault1}

真机串口输出 \code{BPH0abde} 后不再有动静。对照表~\ref{tab:boot-marks}，\code{d} 表示进入了第一次 \code{kfree()}，\code{e} 表示 \code{memset} 已完成，缺少的 \code{f} 表示 \code{acquire(\&kmem.lock)} 没有返回。

当时只有 CPU0 在运行，不可能有别的核持锁。真正的原因在于 \code{acquire()} 的实现：

\begin{asmcode}
retry:  ldaxr w1, [x0]      // 独占加载，建立“独占监视”
        cbnz  w1, retry
        stxr  w2, w3, [x0]  // 独占存储，失败时 w2 != 0
        cbnz  w2, retry
\end{asmcode}

\code{ldaxr} 读到锁为 0，但 \code{stxr} 一直返回失败——这是活锁，不是死锁。当时普通 RAM 被映射成 \emph{Normal Non-cacheable}，也没有 Inner Shareable 属性，数据缓存也没有打开。在 Cortex-A53 上，独占监视器依赖缓存一致性机制工作，对这种属性的内存，独占存储可能永远不成功。

修复包括四处：增加 \code{MT\_NORMAL}（Write-Back）类型并让 RAM 使用它；页表项加 \code{PTE\_SH\_INNER}；\reg{TCR\_EL1} 的页表遍历属性设为 Write-Back、Inner Shareable；\reg{SCTLR\_EL1} 打开 C 和 I 位。之后串口变成 \code{BPH0abdefc1234}，系统顺利进入 shell。

\subsection{故障二：\code{ls} 跳到了不属于它的代码}

打开缓存之后，在 shell 里执行 \code{ls} 出现用户异常：

\begin{txtcode}
usertrap(): unexpected esr=0x92000047 ec=0x24 pid=3
            elr=0x13dc far=0x13f48
\end{txtcode}

EC$=$\code{0x24} 是来自 EL0 的数据中止，ISS 表示 L3 翻译错误、写访问。奇怪的是 \code{ls} 的代码只到 \code{0xbd7}，PC 却在 \code{0x13dc}。

原因在于 D-cache 与 I-cache 不一致。\code{exec()} 通过内核高地址把新程序的指令写进物理页，新指令此时还在 D-cache 里；用户态从低地址取指，走的是 I-cache，里面可能还留着同一虚拟地址上一个程序（\code{sh}）的指令。AArch64 不保证数据写入对取指自动可见，必须显式维护：

\begin{ccode}
void uvmsync_icache(pagetable_t pagetable, uint64 sz) {
  for(uint64 va = 0; va < PGROUNDUP(sz); va += PGSIZE){
    uint64 kva = uva2ka(pagetable, va);
    if(kva == 0) continue;
    for(uint64 p = kva; p < kva + PGSIZE; p += 64)
      asm volatile("dc cvau, %0" :: "r"(p) : "memory");   // 数据清到统一点
  }
  asm volatile("dsb ish" ::: "memory");
  asm volatile("ic ialluis\n\tdsb ish\n\tisb" ::: "memory"); // 广播作废所有核的 I-cache
}
\end{ccode}

第一版用 \code{ic ivau} 按内核高地址作废 I-cache，仍然失败：同一物理页有内核高地址和用户低地址两个别名，按其中一个别名维护不能保证覆盖另一个。改用 \code{ic iallu} 作废整个 I-cache 后问题消失。但打开四核以后它又以另一种形式回来了：\code{ic iallu} 只作废\emph{执行它的那个核}的 I-cache，\code{exec()} 在一个核上装好新程序，进程迁移到另一个核后，那里仍能命中旧程序的指令——表现为 \code{ELR} 指向一条根本不访存的指令却报数据中止（第~\ref{chap:drivers} 章摄像头一节）。最终版本用广播到 Inner Shareable 域的 \code{ic ialluis}。代价是每次 \code{exec()}/\code{fork()} 清一次 I-cache，对 xv6 的规模完全可以接受。\code{exec()}、\code{uvminit()}（第一个进程）、\code{uvmcopy()}（\code{fork}）三处都会调用它。

\subsection{故障三：\code{sbrk} 之后的旧 TLB}

进一步的日志显示，上面那次异常还有另一层原因：\code{sh} 通过 \code{sbrk()} 新增了 \code{0x13000} 这一页，但 \code{growproc()} 修改页表后没有作废 TLB，CPU 继续使用之前缓存的“该地址无效”的翻译结果。修复是在页表扩展或收缩后执行广播式 TLB 失效：

\begin{asmcode}
        dsb  ishst          // 确保页表写入完成
        tlbi vmalle1is      // 作废 Inner Shareable 域内所有核的 EL1 TLB
        dsb  ish
        isb
\end{asmcode}

\code{vmalle1is} 中的 \code{is} 很关键：多核以后，同一进程的线程可能正在别的核上运行，只作废本核 TLB 是不够的。收缩地址空间时顺序必须是“清页表项 → 广播作废 TLB → 释放物理页”，否则别的核可能通过陈旧的 TLB 访问一页已经被重新分配的内存。

\begin{note}
  经验：在 AArch64 上，改了“将要执行的内存”要做 D/I-cache 同步，改了“页表”要做 TLB 失效，两者不能互相替代。QEMU 两件事都不检查，所以必须上真机验证。
\end{note}

\section{实验分支：引用计数、写时复制与页缓存}
\label{sec:sharedpages}

前面各节描述的都是当前源码树。本节换一个角度，记录实验分支 \code{wip-cow-mmap-ldso} 的设计，它没有整体合入当前源码树：其中的写时复制 fork 已经以另一种方式实现在主线里（第~\ref{sec:cow-fork} 节），区别在于主线用 \code{struct page} 的 \code{refcount} 而不是单独的 \code{pageref[]} 数组，\code{PTE\_COW} 用第 57 位（第 55 位已被 \code{PTE\_SHARED} 占用），用每个地址空间的 \code{vm->lock} 而不是全局的 \code{pte\_lock}；页缓存仍只在实验分支里，\code{mmap} 的实现见第~\ref{sec:mmap} 节。下面的设计是让多个地址空间共用物理页——fork 后的父子进程、所有进程里的 \code{libc.so} 代码、同一文件的多个 \code{mmap}——所需的基础，前提是先回答“这一页还有谁在用”。

\subsection{物理页引用计数}

实验分支写于伙伴系统之前，当时 \file{kalloc.c} 还是空闲链表，没有页描述符，只好为它能分配的每一页另开一个 16 位计数数组，以物理页号为下标，受 \code{kmem.lock} 保护。当前源码树的 \code{struct page} 已经带有 \code{refcount} 和 \code{get\_page()}/\code{put\_page()}（第~\ref{sec:buddy} 节），把这套设计移植过来时，下面的 \code{pageref[]} 就直接换成 \code{mem\_map[pfn].refcount}：

\begin{ccode}
static uint16 pageref[PHYSTOP / PGSIZE];

void *kalloc(void)      // 取出空闲页时置为 1
void  kref_inc(uint64 pa)   // 再多一个映射（COW、页缓存、共享库）
int   kref_get(uint64 pa)
void  kfree(void *va)   // 减一；减到 0 才真正放回空闲链表
\end{ccode}

于是“释放”的含义变成了“放弃我这一份”：\code{uvmunmap()}、\code{growproc()} 缩小内存、进程退出，一律调用 \code{kfree()}，共享页只会在最后一个使用者离开时回到空闲链表。不归 \code{kalloc} 管的物理页（相机 DMA 缓冲、模块区）不计数，也不允许传给这些函数。

这一步也修正了第一版共享 libc 时的隐患：那时共享页用软件位 \code{PTE\_SHARED} 标记、各处“见到就别释放”，结果 \code{sbrk(-sz)} 走的是 \code{growproc()} 自己的释放循环，绕过了检查，把所有进程共用的 libc 代码页还给了空闲链表。usertests 的 \code{sbrkbugs} 正是这样暴露问题的：子进程释放全部内存后，父进程、\code{sh}、\code{init} 相继在 libc 的同一地址遇到未定义指令（EC $=0$），最终 \code{panic: init exiting}。有了引用计数，这类特判全部消失。

\subsection{写时复制}
\label{sec:cow}

fork 之后父子共享的可写页在两边都是“只读 + \code{PTE\_COW}”。\code{PTE\_COW} 是 ARMv8 留给软件的描述符位 55（MMU 忽略 \code{[58:55]}），它的意思是“只读只是因为共享，写了就复制”。写这种页会产生权限错误：

\begin{ccode}
} else if(ec == 0x24 && ((esr >> 6) & 1) && (esr & 0x3c) == 0x0c &&
          (cow = uvm_cowfault(p->vm->pagetable, r_far_el1())) != -1){
  if(cow == -2){                       // 没有空闲页可复制
    printf("pid %d %s: out of memory for a copy-on-write fault at %p\n", ...);
    p->killed = 1;
  }
  // 否则返回用户态，那条 str 重新执行，这次写进自己的页
}
\end{ccode}

判断条件逐位对应 \reg{ESR\_EL1}：EC \code{0x24} 是来自 EL0 的数据中止，ISS 的 bit 6（WnR）为 1 表示写，DFSC \code{0b0011xx} 是权限错误。\code{uvm\_cowfault()} 在 \code{pte\_lock} 下处理这一页：

\begin{ccode}
attrs = (*pte & ~(PA2PTE(~0ULL) | PTE_COW | PTE_AP(3))) | PTE_U;  // EL0 可写
if(kref_get(pa) == 1){              // 别人都已复制或退出：直接收归己有
  *pte = PA2PTE(pa) | attrs;
} else {
  mem = kalloc();                   // 复制一份
  memmove(mem, P2V(pa), PGSIZE);
  if((attrs & PTE_UXN) == 0)        // 可执行页还要同步 I-cache
    uvmsync_icache_kva(mem);
  *pte = PA2PTE(V2P(mem)) | attrs;
  kfree(P2V(pa));                   // 放弃对共享页的那一份引用
}
\end{ccode}

改完页表项后 \code{flush\_tlb()} 广播作废，旧的只读翻译从所有 CPU 上消失。真正“只读”的页（共享库代码、\code{PROT\_READ} 映射）没有 \code{PTE\_COW}，写它们仍是致命错误。

\paragraph{\code{copyout()} 也要遵守.} 内核写用户内存走的是自己的直接映射，用户页表项的 AP 位对它不起作用。如果不加检查，子进程 \code{read(fd, buf, n)} 时内核会把数据直接写进仍与父进程共享的物理页，父进程的内存就被改了；同样，\code{read(fd, libc\_text, n)} 能让一个进程改写所有进程共用的 libc 代码。所以 \code{copyout()} 对每一页先在 \code{pte\_lock} 下检查：COW 页先复制，只读页拒绝（返回 $-1$），并且\emph{持锁完成这一页的拷贝}，免得另一个线程的 \code{fork()} 恰好在检查和写入之间把它重新变成共享页。

\paragraph{一把锁.} \code{pte\_lock} 串行化所有“把页变成 COW”和“把 COW 页变回可写”的操作：fork 标记、写保护异常、\code{copyout()}、区间撤销映射。没有它，一个线程可能正把 \code{kref\_get()==1} 的页就地改成可写，而同一地址空间里另一个线程的 \code{fork()} 恰好把同一页共享给了子进程。

\paragraph{效果与代价.} 新增的 \code{cowtest} 实测：8\,MiB 的进程 fork 只用 52\,KiB；子进程写遍 8\,MiB 后才真正复制 8\,MiB；连续 fork 50 次空闲内存不变。代价是内存从“fork 时不够就失败”变成“写的时候不够才失败”：usertests 的 \code{countfree()} 子进程先用 \code{sbrk} 吃光内存，再写一个仍与父进程共享的数据页，于是在写保护异常里因没有空闲页而被杀——这正是 Linux OOM 的行为。另一个副作用是 fork 太便宜，进程数不再被内存限制住；动态增长的进程表和线性扫描的调度器经不起几万个进程，因此 \file{param.h} 加了 \code{NPROC\_MAX} $=1024$ 的上限，作用相当于 Linux 的 \code{threads-max}。

\subsection{页缓存}

要让不同进程共享同一个\emph{文件}的页（典型是 \code{libc.so} 的代码），内核需要一个以“文件 + 页号”为键的缓存。\file{kernel/pcache.c} 是一个很小的版本：

\begin{ccode}
struct pcfile {
  struct inode *ip;     // 持有 inode 引用，像 Linux 的 address_space
  uint64 *pa;           // 一页数组：文件第 n 页 -> 物理页
  uint64 lastuse;
};
uint64 pcache_getpage(struct inode *ip, uint64 pgidx);  // 返回 pa，并为调用者 kref_inc
void   pcache_invalidate(struct inode *ip);             // 文件变了：丢弃缓存页
\end{ccode}

\begin{itemize}
  \item 缓存页是普通的 \code{kalloc()} 页，缓存自己持有一份引用，每个映射再各持一份；文件末尾之后的字节读作 0。读入后统一做一次 I-cache 同步，因为它可能被映射成可执行的库代码。
  \item 缓存最多 8 个文件（\code{NINODE} 只有 50），每个文件最多 512 页（2\,MiB）；满了就按最近使用淘汰。被淘汰文件的 inode 引用不能在持锁时释放，先放进待释放列表，由 \code{pcache\_reap()} 在锁外、自开日志事务再 \code{iput()}。
  \item 一致性：\code{filewrite}、\code{ftruncate} 和 \code{open(O\_TRUNC)} 之后调用 \code{pcache\_invalidate()}。已有的映射保留它们拿到的页，相当于一份 \code{MAP\_PRIVATE} 快照；新的 \code{mmap} 看到新内容。
  \item 加锁顺序固定为“先 \code{ilock(ip)}，再页缓存的睡眠锁”；\code{pcache\_invalidate()} 在 inode 解锁之后调用。
\end{itemize}

\subsection{实验分支里的 \code{mmap()}}

用户地址空间在 \code{sz} 之上多了一块区域 \code{[MMAP\_BASE, MMAP\_TOP)}，即 64\,GiB 到相机窗口下方。每个 \code{struct vmspace} 有一张 32 项的 VMA 表记录这里的区间：

\begin{ccode}
struct vma { uint64 start, end; int prot, flags; };   // start == end 表示空闲
\end{ccode}

\file{kernel/mmap.c} 在调用时就把整段映射好（没有按需调页），所以 VMA 只需记录区间：

\begin{table}[htbp]
  \centering
  \caption{\code{mmap()} 支持的形式}
  \label{tab:mmap}
  \small
  \begin{tabularx}{\linewidth}{@{}lX@{}}
    \toprule
    形式 & 实现 \\
    \midrule
    \code{MAP\_ANONYMOUS} & 新分配的清零页，\code{PROT\_WRITE} 时可写 \\
    文件 + \code{MAP\_PRIVATE} & 页缓存页；只读映射直接共享，可写映射以“只读 + \code{PTE\_COW}”共享，写时复制，缓存中的原页不变 \\
    文件 + \code{MAP\_SHARED} & 只允许只读（没有回写） \\
    \code{MAP\_FIXED} & 先撤销该区间原有的映射，再映射到指定地址 \\
    \code{munmap()} & 可以只撤销区间的一部分，必要时把一个 VMA 切成两个 \\
    \bottomrule
  \end{tabularx}
\end{table}

\code{PROT\_EXEC} 决定 UXN 位；内核永远不执行用户页（PXN）；\code{PROT\_NONE} 映射成 AP $=00$，EL0 无法访问。fork 时 \code{vma\_copy()} 对每个区间调用 \code{uvmcopy\_range()}，因此 mmap 区域同样写时复制；进程退出和 \code{exec()} 替换镜像时 \code{vma\_unmap\_all()} 撤销全部区间。\code{sbrk} 不允许越过 \code{MMAP\_BASE}；相机的 \code{/dev/video0} 仍由驱动自己提供页，\code{sys\_mmap()} 遇到字符设备就交给 \code{chrdev\_mmap()}。

实验分支的 \code{mmaptest} 覆盖了这些路径：只读映射的内容与文件一致；私有可写映射被写后，另一个映射和文件本身都不变；写 \code{PROT\_READ} 映射的子进程被杀；fork 后子进程的写入不影响父进程；匿名页为零；\code{munmap} 掉中间一页后访问它会出错；\code{MAP\_FIXED} 覆盖已有映射；文件被改写后新的映射能看到新内容。下一章的用户态动态链接器 \code{ld.so} 正是用这套接口装载 \code{libc.so} 的（第~\ref{sec:ldso} 节）。

有了页表和物理页分配器，内核就可以为每个进程准备一个独立的地址空间。下一章看进程本身：\code{struct proc} 里有什么，第一个进程怎样造出来，\code{fork()}、\code{exec()} 和调度器怎样使用本章的这些函数。
